\documentclass[11pt]{article}
\usepackage[letterpaper,margin=0.82in]{geometry}
\usepackage[T1]{fontenc}
\usepackage{lmodern}
\usepackage{microtype}
\usepackage{amsmath,amssymb,amsthm,mathtools}
\usepackage{booktabs,longtable,array,tabularx}
\usepackage[dvipsnames,table]{xcolor}
\usepackage{enumitem}
\usepackage{listings}
\usepackage{fancyhdr}
\usepackage{hyperref}
\usepackage{cleveref}
\usepackage{url}
\usepackage{graphicx}
\usepackage{etoolbox}

\hypersetup{colorlinks=true,linkcolor=MidnightBlue,citecolor=OliveGreen,urlcolor=BrickRed,pdfauthor={OpenAI},pdftitle={The beta function in intuitionistic primitive recursive arithmetic}}
\setlist{nosep,leftmargin=1.5em}
\setlength{\parindent}{0pt}
\setlength{\parskip}{0.55em}
\renewcommand{\arraystretch}{1.15}
\pagestyle{fancy}
\fancyhf{}
\lhead{Beta function in iPRA}
\rhead{16 November 2026}
\cfoot{\thepage}
\setlength{\headheight}{14pt}

\newtheorem{theorem}{Theorem}[section]
\newtheorem{lemma}[theorem]{Lemma}
\newtheorem{proposition}[theorem]{Proposition}
\newtheorem{corollary}[theorem]{Corollary}
\theoremstyle{definition}
\newtheorem{definition}[theorem]{Definition}
\newtheorem{remark}[theorem]{Remark}
\newtheorem{warning}[theorem]{Source-audit warning}

\newcommand{\B}{\mathrm B}
\newcommand{\G}{\Gamma}
\newcommand{\N}{\mathbb N}
\newcommand{\Z}{\mathbb Z}
\newcommand{\Q}{\mathbb Q}
\newcommand{\C}{\mathbb C}
\newcommand{\R}{\mathbb R}
\newcommand{\dd}{\,d}
\newcommand{\Rea}{\operatorname{Re}}
\newcommand{\I}{\mathrm I}
\newcommand{\ipra}{\mathrm{iPRA}}
\newcommand{\code}[1]{\texttt{\detokenize{#1}}}
\newcommand{\rising}[2]{(#1)_{#2}}
\newcommand{\statusA}{\textcolor{ForestGreen}{\textbf{A}}}
\newcommand{\statusB}{\textcolor{MidnightBlue}{\textbf{B}}}
\newcommand{\statusC}{\textcolor{BurntOrange}{\textbf{C}}}
\newcommand{\statusD}{\textcolor{BrickRed}{\textbf{D}}}
\newcommand{\abs}[1]{\left|#1\right|}
\newcommand{\norm}[1]{\left\lVert#1\right\rVert}
\newcommand{\ceil}[1]{\left\lceil#1\right\rceil}
\newcommand{\floor}[1]{\left\lfloor#1\right\rfloor}

\lstset{basicstyle=\ttfamily\small,columns=fullflexible,breaklines=true,frame=single,backgroundcolor=\color{gray!5},keywordstyle=\color{MidnightBlue},commentstyle=\color{ForestGreen}}

\title{\textbf{The beta function in intuitionistic primitive recursive arithmetic}\\
\large Effective proofs, source audit, and a Lean/ConciseIPRA formalization route}
\author{Prepared for Vladimir Reshetnikov}
\date{16 November 2026}

\begin{document}
\maketitle

\begin{abstract}
This report reconstructs, in a finitary constructive form suitable for intuitionistic primitive recursive arithmetic, the mathematical claims appearing on the Wikipedia and MathWorld pages for Euler's beta function.  The basic object is not an abstract real-valued function.  It is a primitive-recursive rational approximation algorithm, parameterized by fixed primitive-recursive names for its inputs and accompanied by explicit endpoint, mesh, denominator, and tail certificates.  Within that reading we prove the effective beta integral, symmetry, positivity, Pascal and shift recurrences, integer and binomial values, the beta--gamma identity, transformed integrals, parameter derivatives, reflection and duplication consequences, infinite products, multiplication identities, both Gosper product formulas, incomplete and regularized beta identities, the binomial-CDF relation, multivariate beta recurrences, and quantitative asymptotics.  We also isolate claims requiring additional complex-analysis infrastructure and identify source formulas that need qualification or correction.  A Lean outline specifies the exact definitions, induction invariants, module order, and target formulas for completing a kernel-checked development in ProveIt's \code{Prf PRAAx} and transporting it through \code{pra_iff_native}.
\end{abstract}

\begin{center}
\fbox{\begin{minipage}{0.94\textwidth}
\textbf{Proof-status convention.}  \statusA{} means a complete finite proof is supplied here and its reduction to quantifier-free arithmetic induction is explicit.  \statusB{} means the proof is complete modulo the effective Gamma/elementary-function contracts already developed in ProveIt's Gamma research packet.  \statusC{} means a mathematically precise constructive route and Lean obligations are supplied, but a substantial analytic certificate remains to be built.  \statusD{} marks a false-as-written, branch-ambiguous, or out-of-scope source claim.  The report is a research packet, not a claim that all analytic obligations have already been checked by Lean.
\end{minipage}}
\end{center}

\tableofcontents
\newpage

\section{Scope, theorem inventory, and exact meaning of the claim}

\subsection{The source pages and the constructive target}

The two requested source pages mix several logically different kinds of statement:

\begin{enumerate}
\item identities of finite algebra, such as Pascal's relation and the factorial values;
\item identities of proper or improper real integrals;
\item results about named transcendental functions, such as Gamma, digamma, sine, and exponential;
\item limiting formulas, products, asymptotics, and continued fractions;
\item complex analytic continuation and contour integrals;
\item applications and software remarks that are not themselves mathematical theorems about Beta.
\end{enumerate}

A literal function $\R^2\to\R$ is not an object of first-order iPRA.  The constructive target is instead a family of arithmetical theorems about rational approximation programs.  A real input is a fixed primitive-recursive name, not a quantified arbitrary program.  A domain condition is positive data: lower bounds for real parts, upper bounds for moduli, and lower bounds for denominators that will be inverted.  An equality $F=G$ means that for every requested precision $p$, explicitly chosen finite calculations for $F$ and $G$ differ by at most $2^{-p}$.  This is the same discipline used by ProveIt's Gamma packet.

The report therefore distinguishes three levels:

\begin{description}
\item[Mathematical finite proof.]  Every limit is replaced by a displayed primitive-recursive modulus, every improper endpoint by a displayed tail bound, every division by a displayed apartness margin, and every differentiation by a finite Taylor remainder.
\item[Lean metatheory outline.]  Rational codes, approximation programs, finite sums, and certificates are represented in Lean.  This layer may use ordinary Lean structures to organize the proof.
\item[Internal iPRA derivation.]  The finite invariant is reified as a formula over ProveIt's \code{PTm}; quantifier-free induction produces a term of type $[]\vdash[\code{PRAAx}]\varphi$.  The checked bridge \code{pra_iff_native} then transfers the result to the native concise calculus.
\end{description}

The first level is developed in detail below.  The second and third are specified down to definitions and target formulas, but the complete analytic Lean development is not claimed to have been executed.

\subsection{Coverage at a glance}

The report proves or reduces to existing effective Gamma contracts the following large families.

\begin{longtable}{@{}p{0.27\textwidth}p{0.56\textwidth}p{0.08\textwidth}@{}}
\toprule
Family & Results covered & Status\\
\midrule\endhead
Definition and convergence & Effective integral on $[0,1]$; endpoint tails; primitive-recursive mesh; independence of certificates & \statusA\\
Elementary laws & Symmetry, conjugation, positivity, bounds, Pascal identity, both shift recurrences, rising-factorial shifts & \statusA\\
Discrete values & $B(a,1)$, $B(a,n+1)$, factorial values, binomial relation, $h(a,b)$ coefficient formula and PDE & \statusA\\
Gamma relation & $B(a,b)=\Gamma(a)\Gamma(b)/\Gamma(a+b)$ through a one-dimensional effective limit & \statusB\\
Integral transformations & Trigonometric, infinite-interval, power substitution, fractional-linear substitution, central radical integral & \statusA/\statusB\\
Series and products & General binomial series; effective Euler product; source audit of the displayed symmetric series & \statusA/\statusD\\
Differentiation & All mixed parameter derivatives as logarithmic moments; digamma/polygamma formulas & \statusA/\statusB\\
Special identities & Reflection, product-reflection, duplication, associativity, special values & \statusB\\
Multiplication & MathWorld's general $n$-fold identity; odd and even Gosper products and the $3$- and $4$-fold cases & \statusB\\
Incomplete Beta & Definition, transformed integrals, polynomial values, recurrences, derivative and antiderivative, hypergeometric series & \statusA/\statusB\\
Regularized Beta & Complement, endpoint values, recurrences, binomial-CDF identity & \statusA\\
Continued fraction & Finite continuants and an exact formalization target; convergence/alternating bounds delegated to a positive-remainder lemma & \statusC\\
Multivariate Beta & Gamma quotient, coordinate shifts, Pascal identity, derivatives, fixed-dimensional simplex reduction & \statusA/\statusB\\
Asymptotics & Two-parameter Stirling formula with explicit error; fixed-parameter natural-grid limit & \statusB\\
Complex continuation & Meromorphic Gamma quotient; Pochhammer contour and general-frequency reciprocal integrals outlined & \statusB/\statusC\\
\bottomrule
\end{longtable}

A machine-readable, formula-by-formula ledger is included as \code{evidence/theorem_inventory.csv}.

\subsection{What is reused from ProveIt}

The current repository separates checked arithmetic from research packets.  The ordinary object theory is \code{PRAAx}: intuitionistic first-order logic with equality, a symbol and defining equations for each primitive-recursive function, successor axioms, and induction for quantifier-free formulas.  The \code{Derivations} library already supplies checked semiring, order, finite-sum, finite-product, power, factorial, division, pairing, bounded-search, and related infrastructure.  The \code{ConciseIPRA} library supplies the bi-interpretation and the theorem \code{pra_iff_native}.

The Gamma packet supplies a mathematically detailed but not fully internalized effective construction of Gamma, elementary functions, finite Riemann integration, reflection, multiplication, derivative certificates, and Stirling estimates.  Its own documentation explicitly says that the full analytic Lean development and the production of \code{Prf PRAAx} derivations remain work.  This report reuses those contracts transparently: every result marked \statusB{} inherits that qualification.

\subsection{Constructive logic policy}

No proof below uses excluded middle for equality of real names, countable choice, an abstract Cauchy-completeness axiom, unrestricted dominated convergence, or an unbounded search for a modulus.  Case distinctions are on natural numbers or rational comparisons.  When a reciprocal is needed, an explicit positive rational lower bound is part of the hypotheses.  When a source page writes ``$z\ne n$,'' the constructive replacement is a rational $\delta>0$ satisfying $|z-n|\ge\delta$ for the finitely many relevant integers.

\section{Finitary analytic foundations inside iPRA}
\label{sec:foundations}

\subsection{Codes and regular names}

Integers are coded by a sign and a natural magnitude.  A rational code is a pair $(u,v)$ with $u\in\Z$ and $v\in\N_{>0}$.  Equality and order are decided by cross multiplication; gcd normalization is optional.  A complex rational is a pair of rational codes.  All of these operations are primitive recursive after denominators are carried as positive natural numbers.

\begin{definition}[Regular names]
A regular real name is a fixed primitive-recursive map $x:\N\to\Q$ with
\[
 |x(p)-x(q)|\le 2^{-p}+2^{-q}.
\]
A complex name is a pair of regular real names.  We write $x\simeq y$ if
\[
 \forall p\quad |x(p+3)-y(p+3)|\le 2^{-p}.
\]
A positive margin for $x$ is a rational $d>0$ together with a proof that every sufficiently accurate approximation is at least $d$; for a complex name it is a lower bound on the squared modulus or on the $1$-norm.
\end{definition}

The fixed shift $3$ is inessential.  It absorbs the two regularity errors and makes transitivity immediate after refinement.  The definition does not require a sort containing all primitive-recursive programs.  Each theorem is a schema for the particular program names occurring in it.

\subsection{Arithmetic error lemmas}

Every analytic construction below ultimately uses the following finite identities.  They are especially suitable for ProveIt's embedded arithmetic library because their proofs are inductions on a finite length.

\begin{lemma}[Finite error budgets]
Suppose $|a_j-b_j|\le\varepsilon$ for $j<s$.
\begin{align*}
 \left|\sum_{j<s}a_j-\sum_{j<s}b_j\right|&\le s\varepsilon,\\
 \left|\prod_{j<s}a_j-\prod_{j<s}b_j\right|&\le sM^{s-1}\varepsilon
 \quad\text{if }|a_j|,|b_j|\le M,\\
 |a^{-1}-b^{-1}|&\le \varepsilon/d^2
 \quad\text{if } |a|,|b|\ge d.
\end{align*}
If only $|a|\ge d$ is known and $|a-b|\le d/2$, the last bound may be replaced by $2\varepsilon/d^2$.
\end{lemma}

\begin{proof}
The sum estimate is the triangle inequality.  For products use the exact telescoping identity
\[
 \prod_{j<s}a_j-\prod_{j<s}b_j
 =\sum_{r<s}(a_r-b_r)\prod_{j<r}a_j\prod_{r<j<s}b_j.
\]
The inverse identity is $(a^{-1}-b^{-1})=(b-a)/(ab)$.  After rational values replace every absolute-value comparison by finitely many integer inequalities.  The telescope is proved by quantifier-free induction on $s$.
\end{proof}

A deliberately wasteful precision shift is often simpler than binary logarithms.  If $C\ge1$ is an integer, then
\[
 C2^{-(C+p+3)}<2^{-p}.
\]
Thus every displayed integer majorant can be converted directly into a primitive-recursive precision allocation.

\subsection{Positive powers, exponential, and logarithm}

For a positive real name $t$ with rational margins $0<d\le t\le D$, define $\log t$ by the atanh series
\[
 \log t=2\sum_{j\ge0}\frac{v^{2j+1}}{2j+1},\qquad v=\frac{t-1}{t+1}.
\]
If $D^{-1}\le t\le D$, then $|v|\le1-1/D$ and a primitive-recursive geometric cutoff gives any requested precision.  Define $\exp w$ by its factorial series.  On $|w|\le R$, once $m+1\ge2R$ the tail is bounded by
\[
 2\frac{R^{m+1}}{(m+1)!}.
\]
Finally, for positive $t$ and a complex name $z$, set
\[
 t^z:=\exp(z\log t).
\]
No complex logarithm or branch choice is involved because the base is positive real.  Finite Cauchy products prove $\exp(u+v)=\exp u\exp v$ and hence the usual power laws on certified compact domains.

\subsection{Finite integration}

Let $f$ be supplied on a rational interval $[c,d]$ with a rational Lipschitz majorant $L$.  The left Riemann sum on $m$ equal subintervals differs from every refinement, and hence from the diagonal name defined by those refinements, by at most
\[
 \frac{L(d-c)^2}{m}.
\]
Choose $m\ge 2^{p+4}(1+\ceil{L(d-c)^2})$ and evaluate every summand with a total arithmetic error below $2^{-p-4}$.  This defines the integral as a regular primitive-recursive name.  There is no appeal to a general completeness theorem.

\begin{lemma}[Finite fundamental theorem]
Suppose a fixed primitive-recursive approximation to $F$ satisfies, on the certified interval,
\[
 |F(t+h)-F(t)-hf(t)|\le C h^2.
\]
Then the Riemann-sum name for $\int_c^d f(t)\,dt$ equals $F(d)-F(c)$.
\end{lemma}

\begin{proof}
Sum the displayed relation over the mesh.  The main terms telescope exactly; the total remainder is at most $Cm h^2=C(d-c)^2/m$.  Choose the mesh from the requested output precision.  This is a finite sum identity followed by quantifier-free induction.
\end{proof}

The same argument proves integration by parts and substitution whenever the derivative and quadratic remainder of the substitution are supplied.  Improper integrals are not primitive notions: first truncate, then add an explicit tail certificate.

\subsection{Differentiability and analyticity}

A derivative certificate for $F$ at $z$ is a program for $F'(z)$ and a rational majorant $C$ such that
\[
 |F(z+h)-F(z)-hF'(z)|\le C|h|^2
\]
on a certified neighborhood.  Higher derivatives use a finite polynomial with an explicitly bounded remainder.  A holomorphic-disk certificate consists of finite coefficient approximations, a larger rational radius, and a geometric coefficient majorant.  These are positive data, not predicates decided about arbitrary functions.

\begin{proposition}[Finite-precision interpretation]
Any proof in this report that supplies endpoint tails, compact derivative bounds, finite algebraic identities, and denominator margins translates, for each fixed tuple of primitive-recursive names and each precision $p$, to a finite rational statement.  Its loop invariants are equations or inequalities between primitive-recursive terms and are therefore eligible for quantifier-free induction in iPRA.
\end{proposition}

\begin{proof}
Fix $p$ before beginning the proof.  Replace each named value by a rational approximation at the allocated precision.  Replace every integral by a finite sum, every infinite sum or product by its displayed cutoff, and every reciprocal by cleared multiplication using its margin.  All remaining calculations are bounded loops and integer comparisons.  Universal introduction over $p$ is intuitionistically valid after the finite derivation has been constructed.
\end{proof}

\section{The effective beta integral}
\label{sec:definition}

\subsection{Domain data and endpoint tails}

Fix complex names $a,b$ and integers $A,B,R\ge1$ with certificates
\[
 \Rea a\ge\alpha:=1/A,
 \qquad \Rea b\ge\beta:=1/B,
 \qquad |a|,|b|\le R.
\]
For $0<t<1$, put
\[
 f_{a,b}(t)=t^{a-1}(1-t)^{b-1},
\]
using only the real logarithms of the positive bases $t$ and $1-t$.
For $0<\varepsilon\le1/2$, we have
\begin{align}
 \int_0^\varepsilon |f_{a,b}(t)|\,dt
 &\le \frac{2\varepsilon^\alpha}{\alpha},
 \label{eq:lefttail}\\
 \int_{1-\varepsilon}^1 |f_{a,b}(t)|\,dt
 &\le \frac{2\varepsilon^\beta}{\beta}.
 \label{eq:righttail}
\end{align}
Indeed, on $[0,1/2]$ the factor $(1-t)^{\Rea b-1}$ is at most $2$, because its exponent is greater than $-1$; the other endpoint is symmetric.

For requested precision $p$, let
\[
 L=p+A+B+10,
 \qquad \varepsilon=2^{-ABL}.
\]
Then both powers $\varepsilon^{1/A}$ and $\varepsilon^{1/B}$ are at most $2^{-L}$, so the sum of the two tails is below $2^{-p-5}$.  This cutoff is intentionally enormous but primitive recursive and easy to verify.

\subsection{A compact derivative bound}

On $[\varepsilon,1-\varepsilon]$,
\[
 f'_{a,b}(t)=f_{a,b}(t)
 \left(\frac{a-1}{t}-\frac{b-1}{1-t}\right).
\]
A safe rational majorant is
\[
 |f_{a,b}(t)|\le \varepsilon^{-2(R+1)},
 \qquad
 |f'_{a,b}(t)|\le C:=2(R+1)\varepsilon^{-2R-3}.
\]
Choose
\[
 m=2^{p+7}(1+\ceil C).
\]
The compact Riemann error is then below $2^{-p-7}$.  Evaluate each power in the sum so that the total arithmetic error is also below $2^{-p-7}$.  Combining this with \eqref{eq:lefttail}--\eqref{eq:righttail} gives a rational approximation within $2^{-p-3}$.

\begin{definition}[Effective beta name]
The program $\operatorname{BetaApprox}(a,b,A,B,R,p)$ is the above truncated left Riemann sum, evaluated with the fixed elementary-function precision budget.  The regular name it determines is denoted
\[
 \B(a,b)=\int_0^1 t^{a-1}(1-t)^{b-1}\,dt.
\]
\end{definition}

\begin{theorem}[Existence, regularity, and certificate independence]\label{thm:beta-name}
The effective beta approximation is a primitive-recursive regular name.  If larger valid values of $A,B,R$ are used, the resulting names are equal.  Thus the notation $\B(a,b)$ is independent of the supplied valid domain certificate.
\end{theorem}

\begin{proof}
Every cutoff, mesh length, rational arithmetic precision, and loop body is primitive recursive.  Compare two outputs using a common refinement whose endpoint cutoff is the smaller of the two endpoint intervals and whose mesh refines both meshes.  Each output differs from the common sum by the sum of its displayed endpoint, mesh, and arithmetic budgets.  At requested precision $p+4$ this total is below $2^{-p}$.  The common-refinement statement is a finite identity about sums and is proved by induction on the mesh length.
\end{proof}

\subsection{Symmetry, conjugation, positivity, and coarse bounds}

\begin{theorem}[Basic laws]\label{thm:basic-laws}
On the certified right-half-plane domain:
\begin{align*}
 \B(a,b)&=\B(b,a),\\
 \overline{\B(a,b)}&=\B(\overline a,\overline b),\\
 |\B(a,b)|&\le \B(\Rea a,\Rea b).
\end{align*}
For positive real $a,b$ the value is positive.  If additionally $a,b\le R$, then
\[
 3^{1-2R}\le \B(a,b)\le 2A+2B+2^{2R+3}.
\]
The constants are deliberately coarse.
\end{theorem}

\begin{proof}
Symmetry is the finite substitution $t\mapsto1-t$; on equal meshes it is reversal of a finite list.  Conjugation commutes with the finite elementary approximations, and the modulus inequality is the triangle inequality on each finite sum followed by the endpoint bounds.

For positivity, every summand is nonnegative.  On $[1/3,2/3]$ and $0<a,b\le R$, each factor is at least $3^{1-R}$, so the integral over an interval of length $1/3$ is at least $3^{1-2R}$.  For the upper bound, use the endpoint estimates with $\varepsilon=1/2$ and a crude bound $2^{2R+2}$ on the middle half.  The positive lower bound is an explicit inversion margin for $1/\B(a,b)$.
\end{proof}

\subsection{The direct finite statement internalized in iPRA}

The definition above is external notation for a family of integer-coded programs.  A typical internal target is the formula
\[
 \forall p\quad
 \left|\operatorname{BetaApprox}(a,b,A,B,R,p+4)
      -\operatorname{BetaApprox}(a,b,A',B',R',p+4)
 \right|\le2^{-p},
\]
where the absolute-value inequality is cleared to a comparison of integer products.  The hypotheses saying that both tuples are valid domain certificates are finite Boolean conjunctions.  The proof uses induction only on the Riemann mesh and on the finite elementary-function sums.  It never uses induction on the unbounded statement that two names are equal.

\section{Recurrences, integer values, binomial coefficients, and finite generating identities}
\label{sec:elementary}

\subsection{Pascal and integration-by-parts identities}

The key point is that the most useful beta identities do not require Gamma.

\begin{theorem}[Pascal identity]\label{thm:pascal}
For $\Rea a,\Rea b>0$,
\[
 \boxed{\B(a,b)=\B(a+1,b)+\B(a,b+1).}
\]
\end{theorem}

\begin{proof}
On every truncated interval,
\[
 t^{a-1}(1-t)^{b-1}\bigl(t+(1-t)\bigr)
 =t^a(1-t)^{b-1}+t^{a-1}(1-t)^b.
\]
This is an exact rational identity after the elementary functions have been approximated with a shared error budget.  Sum it over the mesh and then remove the endpoints with \eqref{eq:lefttail}--\eqref{eq:righttail}.
\end{proof}

\begin{lemma}[Balanced integration by parts]\label{lem:balanced-ibp}
For $\Rea a,\Rea b>0$,
\[
 a\B(a,b+1)=b\B(a+1,b).
\]
\end{lemma}

\begin{proof}
The derivative
\[
 \frac{d}{dt}\bigl(t^a(1-t)^b\bigr)
 =a t^{a-1}(1-t)^b-bt^a(1-t)^{b-1}
\]
has integral zero on $[\varepsilon,1-\varepsilon]$ up to the finite fundamental-theorem error.  Its boundary terms have moduli at most $\varepsilon^\alpha+\varepsilon^\beta$, so they vanish with the same primitive-recursive cutoff used in the definition of Beta.
\end{proof}

Since $\Rea(a+b)\ge\alpha+\beta$, the denominator $a+b$ has a supplied positive modulus margin.  Solving the two linear equations gives the standard shifts.

\begin{corollary}[Coordinate recurrences]\label{cor:coordinate-recurrence}
\[
 \boxed{\B(a+1,b)=\frac{a}{a+b}\B(a,b),\qquad
 \B(a,b+1)=\frac{b}{a+b}\B(a,b).}
\]
More generally, for $m,n\in\N$,
\[
 \boxed{\B(a+m,b+n)=
 \frac{(a)_m(b)_n}{(a+b)_{m+n}}\B(a,b),}
\]
where $(z)_r=\prod_{j<r}(z+j)$.
\end{corollary}

\begin{proof}
The one-step formulas are obtained from \cref{thm:pascal,lem:balanced-ibp}.  The general formula is induction first on $m$ and then on $n$; the two denominator products concatenate to $(a+b)_{m+n}$.  The invariant is an equality of finite products after denominators are cleared.
\end{proof}

\subsection{Values at natural arguments}

\begin{theorem}[Base and factorial values]\label{thm:integer-values}
For $\Rea a>0$ and $N\in\N$,
\[
 \B(a,1)=\frac1a,
 \qquad
 \boxed{\B(a,N+1)=\frac{N!}{a(a+1)\cdots(a+N)}.}
\]
For positive integers $m,n$,
\[
 \boxed{\B(m,n)=\frac{(m-1)!(n-1)!}{(m+n-1)!}.}
\]
\end{theorem}

\begin{proof}
The first identity follows from the finite fundamental theorem applied to $t^a/a$; the lower endpoint is bounded by $\varepsilon^\alpha/|a|$.  Iterate the second recurrence in \cref{cor:coordinate-recurrence} to obtain the second formula.  Put $a=m$ and simplify the rising product to obtain the third.  Every step is finite induction and factorial arithmetic.
\end{proof}

The alternative form appearing on Wikipedia is
\[
 \B(m,n)=\frac{m+n}{mn}\binom{m+n}{m}^{-1}.
\]
It follows by multiplying numerator and denominator of the factorial expression by $(m+n)/(mn)$.

\begin{corollary}[Binomial coefficient through Beta]\label{cor:binomial-beta}
For $0\le k\le n$,
\[
 \boxed{\binom nk=\frac{1}{(n+1)\B(n-k+1,k+1)}.}
\]
\end{corollary}

\begin{proof}
Substitute $m=n-k+1$ and $r=k+1$ into the factorial formula.  The result is the exact integer identity
$(n+1)\B(n-k+1,k+1)=k!(n-k)!/n!$.
\end{proof}

\subsection{The entire interpolation formula for binomial coefficients}

Define the generalized binomial name by
\[
 \binom nk=\frac{\Gamma(n+1)}{\Gamma(k+1)\Gamma(n-k+1)}.
\]
At a noninteger $k$ with a supplied separation from $0,1,\dots,n$, reflection and recurrence give
\begin{equation}
 \boxed{
 \binom nk=(-1)^n n!\,
 \frac{\sin(\pi k)}{\pi\prod_{j=0}^n(k-j)}.}
 \label{eq:binomial-interpolation}
\end{equation}
A constructively preferable form is the cleared identity
\[
 \pi\binom nk\prod_{j=0}^n(k-j)=(-1)^n n!\sin(\pi k),
\]
which remains meaningful at integer $k$.  At an integer the quotient in \eqref{eq:binomial-interpolation} is $0/0$ and must be replaced by the factorial value, or by a certified limit.  Thus the source's displayed quotient needs this domain qualification.

\subsection{The function $h(a,b)$ and its PDE}

The quotient $(e^a-e^b)/(a-b)$ is inconvenient at $a=b$.  Define instead
\[
 h(a,b)=\int_0^1 e^{ta+(1-t)b}\,dt.
\]
On an apartness domain $|a-b|\ge d>0$, the finite fundamental theorem proves
\[
 h(a,b)=\frac{e^a-e^b}{a-b}.
\]
The integral definition is regular on the diagonal.

\begin{theorem}[Beta values as derivatives of $h$]\label{thm:h-derivatives}
For $m,n\in\N$,
\[
 \boxed{
 \frac{\partial^{m+n}h}{\partial a^m\partial b^n}(0,0)
 =\B(m+1,n+1).}
\]
Moreover,
\[
 \boxed{h=h_a+h_b.}
\]
\end{theorem}

\begin{proof}
Finite Taylor expansion of the exponential, with its factorial remainder, allows termwise parameter differentiation on the compact interval.  At $(0,0)$ the derivative is
\[
 \int_0^1 t^m(1-t)^n\,dt=\B(m+1,n+1).
\]
For the PDE, differentiate the same finite approximation and use $t+(1-t)=1$.  Equivalently, the power-series identity
\[
 \frac{e^a-e^b}{a-b}
 =\sum_{m,n\ge0}\frac{a^m b^n}{(m+n+1)!}
\]
has coefficient $1/(m+n+1)!$; multiplying by $m!n!$ gives the beta value.  Every coefficient statement up to a requested total degree is finite.
\end{proof}

\subsection{Convolution of truncated powers}

Let $t_+^r$ mean $t^r$ for $t>0$ and $0$ for $t\le0$.  For a positive named $t$ with margin $t\ge d>0$ and $\Rea x,\Rea y>0$,
\begin{equation}
 \boxed{
 (u\mapsto u_+^{x-1})*(u\mapsto u_+^{y-1})(t)
 =\B(x,y)t^{x+y-1}.}
 \label{eq:convolution}
\end{equation}
Indeed the convolution is supported on $0<u<t$; put $u=tv$.  The endpoint estimates are the beta endpoint estimates multiplied by the compact factor $t^{x+y-1}$.  The Wikipedia restriction $x,y\ge1$ avoids singular endpoint notation, but the certified construction works throughout the right half-planes.

\section{The beta--gamma identity by a one-dimensional effective limit}
\label{sec:betagamma}

The familiar textbook proof multiplies two improper Gamma integrals, invokes an iterated integral, and makes a two-dimensional change of variables.  That proof can be made constructive, but it introduces more machinery than necessary.  We instead use only the one-dimensional beta integral, the shift recurrence, and the effective Euler construction of Gamma.

\subsection{An explicit beta-to-gamma limit}

Fix domain data
\[
 \Rea a\ge\alpha>0,
 \qquad |a|\le d,
 \qquad |b-1|\le C,
\]
where $d,C$ are integers after upward rounding.  Let
\[
 H=1/\alpha+d!.
\]
For $T\ge1$ and
\[
 N\ge\max\{2T,2,T^2+2CT\},
\]
the following effective estimate is available:
\begin{equation}
 \left|N^a\B(a,b+N)-\Gamma(a)\right|
 \le \frac{2H(T^2+2CT)}{N}
       +2\,4^{d+1}d!\,e^{-T/4}.
 \label{eq:beta-gamma-limit}
\end{equation}

\begin{proof}
Substitute $t=s/N$ in the beta integral.  Then
\[
 N^a\B(a,b+N)=
 \int_0^N s^{a-1}(1-s/N)^{N+b-1}\,ds.
\]
On $0\le s\le T$ put
\[
 D(s)=(N+b-1)\log(1-s/N)+s.
\]
The finite logarithm inequalities
\[
 0\le-\log(1-u)-u\le u^2,
 \qquad |\log(1-u)|\le2u
 \quad(0\le u\le1/2)
\]
follow from a finite geometric sum plus its positive remainder.  They give
\[
 |D(s)|\le (T^2+2CT)/N\le1.
\]
The factorial series gives $|e^w-1|\le2|w|$ for $|w|\le1$.  Multiplying by $s^{\Rea a-1}e^{-s}$ and using
\[
 \int_0^\infty s^{\Rea a-1}e^{-s}\,ds\le H
\]
proves the first term in \eqref{eq:beta-gamma-limit}.

For $s\ge T$, the real part of $N+b-1$ is at least $N/2$ after increasing the fixed lower cutoff if necessary, and hence
\[
 |(1-s/N)^{N+b-1}|\le e^{-s/2}
 \quad(0\le s<N).
\]
For $s\ge1$, $s^{\Rea a-1}\le s^d$, while the single factorial-series term gives
$s^de^{-s/2}\le4^d d!e^{-s/4}$.  The beta tail and the Gamma tail are each at most $4^{d+1}d!e^{-T/4}$.  All inequalities are first proved on finite intervals; the displayed exponential tail is the cutoff certificate.
\end{proof}

For error $2^{-p}$ one safe choice is
\[
 T=4(p+d^2+2d+8),\qquad
 N=\max\!\left\{2T,2,T^2+2CT,\,
 2^{p+4}(T^2+2CT)\lceil H\rceil\right\}.
\]
The assignment is a closed primitive-recursive function.

\subsection{Identification of the constant}

By the finite recurrence,
\begin{equation}
 \B(a,b+N)=\frac{(b)_N}{(a+b)_N}\B(a,b).
 \label{eq:beta-shift-N}
\end{equation}
The Euler approximants to Gamma give, for $N\ge2$, the exact finite identity
\begin{equation}
 N^a\frac{(b)_N}{(a+b)_N}
 =\left(\frac{N}{N-1}\right)^a
   \frac{\Gamma_{N-1}(a+b)}{\Gamma_{N-1}(b)}.
 \label{eq:euler-ratio-exact}
\end{equation}
Every factor in this equation is a finite product.  The effective convergence of the Gamma approximants, together with denominator margins for Gamma in the right half-plane, therefore gives
\[
 N^a\frac{(b)_N}{(a+b)_N}\longrightarrow
 \frac{\Gamma(a+b)}{\Gamma(b)}
\]
with a primitive-recursive modulus.  Combining this with \eqref{eq:beta-gamma-limit} and \eqref{eq:beta-shift-N} first proves the cleared identity
\[
 \B(a,b)\Gamma(a+b)=\Gamma(a)\Gamma(b).
\]
Gamma is positive on the positive real axis and has an explicit nonzero margin on certified compact right-half-plane domains.  Hence guarded inversion yields the central result.

\begin{theorem}[Beta--Gamma identity]\label{thm:beta-gamma}
For certified $\Rea a,\Rea b>0$,
\[
 \boxed{\B(a,b)=\frac{\Gamma(a)\Gamma(b)}{\Gamma(a+b)}.}
\]
The proof has a primitive-recursive modulus and uses no general Fubini theorem or multidimensional change of variables.
\end{theorem}

\subsection{Immediate functional consequences}

The Gamma recurrence and reflection formula now produce several source-page identities.  Every division below requires the displayed right-half-plane condition or an explicit pole/sine separation.

\begin{corollary}[Associativity identity]\label{cor:beta-assoc}
Whenever all factors are defined,
\[
 \boxed{\B(p,q)\B(p+q,r)=\B(q,r)\B(q+r,p).}
\]
\end{corollary}

\begin{proof}
Both sides are
$\Gamma(p)\Gamma(q)\Gamma(r)/\Gamma(p+q+r)$ after exact cancellation.  A division-free proof multiplies both sides by the common Gamma denominators first.
\end{proof}

\begin{corollary}[Reflection consequences]\label{cor:beta-reflection}
In the integral strip $0<\Rea x<1$,
\[
 \boxed{\B(x,1-x)=\frac{\pi}{\sin(\pi x)}.}
\]
More generally, where the factors have certified margins,
\[
 \boxed{\B(x,y)\B(x+y,1-y)=\frac{\pi}{x\sin(\pi y)}.}
\]
\end{corollary}

\begin{proof}
The first formula is Euler reflection
$\Gamma(x)\Gamma(1-x)=\pi/\sin(\pi x)$.  For the second, Gamma cancellation leaves
$\Gamma(y)\Gamma(1-y)/x$.  Constructively it is preferable to prove the cleared sine identities before inversion.
\end{proof}

At $x=1/2$ this gives
\[
 \B(1/2,1/2)=\pi,
 \qquad \Gamma(1/2)^2=\pi.
\]
The effective Gamma construction is positive on the positive real axis, so the
positive-square-root specification yields the source-page special value
$\Gamma(1/2)=\sqrt\pi$.

\begin{corollary}[Duplication identity]\label{cor:beta-duplication}
For certified arguments,
\[
 \boxed{2^{z-1}\B(z/2,z/2)=\B(1/2,z/2).}
\]
Also
\[
 \boxed{\B(p,p)\B(p+1/2,p+1/2)=
 \frac{\pi}{2^{4p-1}p}.}
\]
\end{corollary}

\begin{proof}
Both formulas follow by substituting the Gamma quotient and applying Legendre duplication.  For the second,
\[
 \frac{[\Gamma(p)\Gamma(p+1/2)]^2}
 {\Gamma(2p)\Gamma(2p+1)}
 =\frac{2^{2-4p}\pi\Gamma(2p)^2}
 {2p\Gamma(2p)^2}.
\]
All cancellations can be carried out after clearing denominators.
\end{proof}

\subsection{Special positive integrals}

For a positive rational $q$ and a certified $0<\Rea a<q$, the substitution $u=t^q$ in the infinite-interval beta representation gives
\begin{equation}
 \boxed{
 \int_0^\infty\frac{t^{a-1}}{1+t^q}\,dt
 =\frac{\pi}{q\sin(\pi a/q)}.}
 \label{eq:p-integral}
\end{equation}
In particular, for real $q>1$,
\[
 \int_0^\infty\frac{dt}{1+t^q}
 =\frac{\pi}{q}\csc(\pi/q),
\]
which is MathWorld's formula.  The lower and upper tails are bounded respectively by $\varepsilon^{\Rea a}/\Rea a$ and $R^{\Rea a-q}/(q-\Rea a)$, so the substitution is finite and effective.

\section{Integral transformations and central-beta formulas}
\label{sec:transformations}

Every substitution in this section is first performed on a closed rational subinterval, where its derivative has an explicit bound.  The removed endpoints are then controlled by the beta tail estimates.  Thus the formulas are consequences of finite telescoping sums, not of an unrestricted change-of-variables theorem.

\subsection{Trigonometric and infinite-interval forms}

\begin{theorem}[Standard transformed integrals]\label{thm:transformed-integrals}
For $\Rea a,\Rea b>0$,
\begin{align}
 \B(a,b)
 &=2\int_0^{\pi/2}
       \sin^{2a-1}\theta\,\cos^{2b-1}\theta\,d\theta,
 \label{eq:trig-beta}\\
 &=\int_0^\infty
       \frac{u^{a-1}}{(1+u)^{a+b}}\,du,
 \label{eq:infinity-beta}\\
 &=q\int_0^1
       t^{qa-1}(1-t^q)^{b-1}\,dt
 \qquad(q>0).
 \label{eq:power-beta}
\end{align}
\end{theorem}

\begin{proof}
For \eqref{eq:trig-beta}, use $t=\sin^2\theta$, so
$dt=2\sin\theta\cos\theta\,d\theta$.  Positivity of sine and cosine on the first quadrant fixes the power convention.  The endpoint error after cutting off $\theta$ is bounded by the corresponding beta endpoint error because $t=\sin^2\theta$ is monotone and has an explicit primitive-recursive inverse obtained by bisection.

For \eqref{eq:infinity-beta}, use $t=u/(1+u)$, so $dt=du/(1+u)^2$.  A lower cutoff $u=\eta$ contributes at most $\eta^\alpha/\alpha$; an upper cutoff $u=T$ contributes at most $T^{-\beta}/\beta$.

For \eqref{eq:power-beta}, use $u=t^q$.  On a truncated interval the derivative is $qt^{q-1}$; the endpoint powers already provide the required tails.  A rational positive name for $q$ includes upper and lower margins, so the bisection evaluator for $t^q$ is primitive recursive.
\end{proof}

MathWorld's formula
\[
 \int_0^{\pi/2}\sin^n x\cos^m x\,dx
 =\frac12\B\!\left(\frac{n+1}{2},\frac{m+1}{2}\right)
\]
is \eqref{eq:trig-beta} with renamed parameters.  For nonnegative integers $m,n$, \cref{thm:integer-values} reduces it to factorial arithmetic.

\subsection{A fractional-linear one-parameter identity}

Let $c<1$ be a real name with a supplied margin $1-c\ge d>0$.  Put
\[
 u=\frac{1-t}{1-ct}.
\]
As $t$ increases from $0$ to $1$, $u$ decreases from $1$ to $0$, and
\[
 1-u=\frac{(1-c)t}{1-ct},
 \qquad
 \left|\frac{du}{dt}\right|=\frac{1-c}{(1-ct)^2}.
\]
Consequently
\begin{equation}
 \boxed{
 \B(a,b)=(1-c)^b
 \int_0^1\frac{(1-t)^{a-1}t^{b-1}}
 {(1-ct)^{a+b}}\,dt.}
 \label{eq:mobius-beta}
\end{equation}
The denominator is at least $d$ on the whole interval.  The Wikipedia page states $c\le1$; at $c=1$ the substitution degenerates and the displayed integrand/prefactor are not a valid ordinary formula.  The correct direct statement is $c<1$, with $c=1$ only as a separately justified limit in parameter ranges where such a limit exists.

\subsection{The central radical integral}

For positive real $z$, start from \eqref{eq:infinity-beta}:
\[
 \B(z,z)=\int_0^\infty\frac{r^{z-1}}{(1+r)^{2z}}\,dr.
\]
In
\[
 I_z=\int_0^{\pi/2}
 \frac{d\theta}
 {(\sin^{1/z}\theta+\cos^{1/z}\theta)^{2z}},
\]
put $s=\tan\theta$ and then $r=s^{1/z}$.  Since
\[
 \sin^{1/z}\theta+\cos^{1/z}\theta
 =(1+s^2)^{-1/(2z)}(1+s^{1/z}),
 \qquad d\theta=\frac{ds}{1+s^2},
\]
we obtain
\[
 I_z=\int_0^\infty\frac{ds}{(1+s^{1/z})^{2z}}
 =z\int_0^\infty\frac{r^{z-1}}{(1+r)^{2z}}\,dr.
\]
Thus
\begin{equation}
 \boxed{
 \B(z,z)=\frac1z\int_0^{\pi/2}
 \frac{d\theta}
 {(\sin^{1/z}\theta+\cos^{1/z}\theta)^{2z}}.}
 \label{eq:central-radical}
\end{equation}
This proof is constructive for positive real $z$ with upper and lower margins.  The source's exclusion $z\ne1$ is unnecessary: at $z=1$ the integrand is $(\sin\theta+\cos\theta)^{-2}$ and the identity remains true.  For complex $z$, the radicals require branch data not stated on the page, so \eqref{eq:central-radical} should not be read branch-free.

\subsection{Gaussian and generalized Gamma integrals}

The effective Gamma scaling substitution $u=ct^q$ gives, for positive real $c,q$ and $\Rea a>0$,
\begin{equation}
 \boxed{
 \int_0^\infty t^{a-1}e^{-ct^q}\,dt
 =\frac{\Gamma(a/q)}{q\,c^{a/q}}.}
 \label{eq:gamma-scaling}
\end{equation}
The special case $q=2$ is the Gaussian integral family used in MathWorld's polar-coordinate derivation.  In iPRA it is simpler to prove \eqref{eq:gamma-scaling} directly in one dimension and then derive the trigonometric beta integral, rather than formalize a two-dimensional polar Jacobian.

\subsection{Finite-dimensional polar identities}

For nonnegative integers $m,n$, the MathWorld calculation can nevertheless be reconstructed without a general measure theory.  Replace the double Gaussian integral over a square $[-T,T]^2$ by a finite rectangular grid.  Partition the first quadrant into rational angular sectors and radial shells.  The discrepancy between the grid and the polar sum is bounded by a Lipschitz constant times the mesh diameter; the Gaussian tail outside radius $T$ has the elementary factorial bound from the Gamma construction.  After the finite angular symmetry factor is extracted, the identity reduces to
\[
 m!n!
 =2(m+n+1)!
 \int_0^{\pi/2}\cos^{2m+1}\theta\sin^{2n+1}\theta\,d\theta.
\]
This route is formalizable, but it is longer than the one-dimensional proof and adds no new beta theorem.  We therefore classify the full polar-grid implementation as \statusC{}, while the resulting identity itself is already \statusA{} by \cref{thm:transformed-integrals,thm:integer-values}.

\section{Series, products, and exact source checking}
\label{sec:series-products}

\subsection{The binomial-series representation}

For $\Rea x,\Rea y>0$, the source pages give
\begin{equation}
 \B(x,y)=\sum_{n=0}^{\infty}
 \frac{(1-x)_n}{(y+n)n!}.
 \label{eq:beta-series}
\end{equation}
Here $(w)_n$ is the rising factorial.

\begin{theorem}[Effective series formula]\label{thm:beta-series}
Equation \eqref{eq:beta-series} holds as an equality of regular names.  On every certified compact domain
\[
 \Rea x\ge\alpha,\quad \Rea y\ge\beta,\quad |x|,|y|\le R,
\]
there is a primitive-recursive cutoff modulus.
\end{theorem}

\begin{proof}
By symmetry write
\[
 \B(x,y)=\int_0^1 t^{y-1}(1-t)^{x-1}\,dt.
\]
Fix the requested precision.  Remove $t\in[1-\delta,1]$ using the endpoint bound $2\delta^\alpha/\alpha$.  On $[0,1-\delta]$ expand
\[
 (1-t)^{x-1}
 =\sum_{n<N}\frac{(1-x)_n}{n!}t^n+R_N(t).
\]
This is the finite Taylor formula for $\exp((x-1)\log(1-t))$.  The elementary exponential remainder and the bound $|\log(1-t)|\le\log(1/\delta)$ give an explicit factorial cutoff $N$.  Integrating the finite polynomial gives
\[
 \sum_{n<N}\frac{(1-x)_n}{n!}\int_0^{1-\delta}t^{y+n-1}\,dt.
\]
Replacing the upper endpoint $1-\delta$ by $1$ contributes another explicitly bounded endpoint term, and
$\int_0^1t^{y+n-1}\,dt=(y+n)^{-1}$.  The finite algebra is exact after guarded inversion because $\Rea(y+n)\ge\beta+n$.  All cutoffs are primitive recursive in the domain data and precision.
\end{proof}

For the incomplete beta function the same proof gives the more useful uniformly geometric series on $0\le z\le1-\delta$:
\[
 \B_z(a,b)=
 \sum_{n=0}^{\infty}
 \frac{(1-b)_n}{(a+n)n!}z^{a+n}.
\]
This will be used in \cref{sec:incomplete}.

\subsection{A displayed symmetric series is false as written}

The Wikipedia page currently displays
\[
 \B(z,z)=2\sum_{n\ge0}
 \frac{(2z+n-1)_n(-1)^n}{(z+n)n!}
 =\lim_{r\to1^-}2\sum_{n\ge0}
 \frac{(-2z)_n r^n}{(z+n)n!},
 \tag{S}
\]
while explicitly declaring $(\cdot)_n$ to be the rising factorial.

\begin{warning}[Counterexample at $z=1$]\label{warn:false-series}
Formula (S) is false with the stated rising-factorial convention.
\end{warning}

\begin{proof}
For $z=1$, the magnitude of the $n$th term in the first series is
\[
 2\frac{(n+1)_n}{(n+1)n!}
 =2\frac{(2n)!}{(n+1)(n!)^2},
\]
which does not tend to zero; the series diverges.  In the second series, $(-2)_n$ vanishes for $n\ge3$, so the Abel limit is
\[
 2\left(1-1+\frac{2}{3\cdot2}\right)
 =\frac23,
\]
not $\B(1,1)=1$.  Thus neither equality can be used in a proof packet without correcting the source formula.
\end{proof}

The exact-rational checker in the archive records this counterexample.  We do not guess the intended correction.

\subsection{The Euler product for Beta}

The source pages give two algebraically equivalent products:
\begin{align}
 \B(x,y)
 &=\frac{x+y}{xy}
 \prod_{k=1}^{\infty}
 \frac{1+(x+y)/k}{(1+x/k)(1+y/k)},
 \label{eq:beta-product-1}\\
 &=\frac{x+y}{xy}
 \prod_{k=1}^{\infty}
 \left(1+\frac{xy}{k(k+x+y)}\right)^{-1}.
 \label{eq:beta-product-2}
\end{align}
The equivalence of their $k$th factors is the exact identity
\[
 \frac{(1+x/k)(1+y/k)}{1+(x+y)/k}
 =1+\frac{xy}{k(k+x+y)}.
\]

\begin{theorem}[Effective infinite product]\label{thm:beta-product}
On a certified compact domain separated from the poles and from $x=0$, $y=0$, and $x+y=0$, the finite products in
\eqref{eq:beta-product-1}--\eqref{eq:beta-product-2} converge to $\B(x,y)$ with a primitive-recursive modulus.
For positive real $x,y$, the tail after $N$ satisfies
\[
 1\le
 \prod_{k>N}\left(1+\frac{xy}{k(k+x+y)}\right)
 \le\exp(xy/N).
\]
\end{theorem}

\begin{proof}
Insert the effective Weierstrass products for $\Gamma(x)$, $\Gamma(y)$, and $\Gamma(x+y)$ into the beta--gamma quotient.  The exponential normalization factors cancel exactly, leaving \eqref{eq:beta-product-1}.  This cancellation occurs already at each common finite cutoff, so only the Gamma product remainder needs to be bounded.

For positive $x,y$, put $u_k=xy/[k(k+x+y)]$.  Then $u_k\ge0$ and
\[
 \sum_{k>N}u_k\le xy\sum_{k>N}\frac1{k^2}\le\frac{xy}{N}.
\]
The finite inequality $\prod(1+u_k)\le\exp(\sum u_k)$ follows by multiplying the finite inequalities $1+u\le e^u$.  The lower bound is immediate.  For complex parameters use $|u_k|\le C/k^2$ once $k$ exceeds the fixed radius, and the finite telescoping product estimate gives a tail bounded by a primitive-recursive multiple of $1/N$.
\end{proof}

For positive $x,y$ this product also gives monotone rational upper and lower approximations to Beta after the prefactor is bounded.  It is therefore an alternative evaluator to the integral name.

\subsection{Finite products before limits}

A useful formalization policy is to prove the identity
\[
 \frac{x+y}{xy}
 \prod_{k=1}^{N}
 \frac{1+(x+y)/k}{(1+x/k)(1+y/k)}
 =
 \frac{x+y}{xy}\,
 \frac{N!\,(x+y+1)_N}{(x+1)_N(y+1)_N}
\]
only after clearing all denominators.  A simpler equivalent form is obtained directly from the Euler Gamma approximants.  The exact choice is not mathematically important; what matters is that the product regrouping is a quantifier-free induction invariant and the analytic content is isolated in an explicit $O(1/N)$ remainder.

\subsection{Products and logarithms}

On positive real domains,
\[
 \log\B(x,y)=\log(x+y)-\log x-\log y
 -\sum_{k\ge1}\log\left(1+\frac{xy}{k(k+x+y)}\right).
\]
Since $0\le\log(1+u)\le u$, the same $xy/N$ bound controls the logarithmic tail.  Finite differentiation of this logarithmic product supplies an independent route to the digamma formulas in \cref{sec:derivatives}; the Gamma quotient is shorter, but the product route may be preferable in an object-theory development because every denominator is visibly separated on a positive compact domain.

\section{Parameter derivatives, digamma formulas, and logarithmic moments}
\label{sec:derivatives}

\subsection{Effective differentiation under the integral}

For $r,s\in\N$, define
\[
 M_{r,s}(a,b)=\int_0^1
 t^{a-1}(1-t)^{b-1}
 (\log t)^r(\log(1-t))^s\,dt.
\]
The logarithms do not destroy effective convergence.  If $0<t\le1/2$ and $\Rea a\ge1/A$, then
\[
 |\log t|^r\le r!(2A)^r t^{-1/(2A)}.
\]
Indeed, with $v=-\log t$, the single term
$e^{v/(2A)}\ge(v/(2A))^r/r!$ proves the inequality.  Also
$|\log(1-t)|\le2t$ on this interval.  Consequently a safe lower-endpoint bound is
\[
 \int_0^\varepsilon
 |t^{a-1}(1-t)^{b-1}
  (\log t)^r(\log(1-t))^s|\,dt
 \le
 4A\,r!(2A)^r\,2^s\,
 \varepsilon^{1/(2A)}.
\]
The upper endpoint has the symmetric bound with $B,s$.  On the compact middle interval, expand
\[
 t^h(1-t)^k
 =\exp\bigl(h\log t+k\log(1-t)\bigr)
\]
to total degree $r+s+1$ with the factorial-series remainder.  The logarithms are bounded by $\log(1/\varepsilon)$, so the remainder constant is explicit.

\begin{theorem}[All mixed parameter derivatives]\label{thm:mixed-derivatives}
For every fixed $r,s\in\N$ and certified right-half-plane inputs,
\[
 \boxed{
 \frac{\partial^{r+s}}{\partial a^r\partial b^s}\B(a,b)
 =M_{r,s}(a,b).}
\]
The derivative means a finite local polynomial with a primitive-recursive remainder bound.
\end{theorem}

\begin{proof}
Choose the endpoint cutoff from the displayed logarithmic tail bounds.  On the compact interval use the finite Taylor expansion in the increments $h,k$.  Integrate the finite polynomial termwise and bound the integrated remainder.  The coefficient of $h^rk^s/(r!s!)$ is precisely $M_{r,s}$.  All degree loops are finite and all precision schedules are primitive recursive.
\end{proof}

In particular,
\[
 \partial_a\B(a,b)=\int_0^1f_{a,b}(t)\log t\,dt,
 \qquad
 \partial_b\B(a,b)=\int_0^1f_{a,b}(t)\log(1-t)\,dt.
\]

\subsection{Digamma and polygamma formulas}

Let
\[
 \psi(z)=\frac{\Gamma'(z)}{\Gamma(z)},
 \qquad
 \psi_1(z)=\psi'(z),
\]
on a certified domain where Gamma has a positive modulus margin.  Differentiate the cleared beta--gamma identity and only then divide.

\begin{theorem}[First and second derivatives]\label{thm:digamma-beta}
\begin{align*}
 \partial_a\B(a,b)
 &=\B(a,b)\bigl(\psi(a)-\psi(a+b)\bigr),\\
 \partial_b\B(a,b)
 &=\B(a,b)\bigl(\psi(b)-\psi(a+b)\bigr),\\
 \partial_b^2\B(a,b)
 &=\B(a,b)\left(
   [\psi(b)-\psi(a+b)]^2+
   \psi_1(b)-\psi_1(a+b)\right),\\
 \partial_a\partial_b\B(a,b)
 &=\B(a,b)\left(
   [\psi(a)-\psi(a+b)]
   [\psi(b)-\psi(a+b)]
   -\psi_1(a+b)\right).
\end{align*}
\end{theorem}

\begin{proof}
The first two identities are the logarithmic derivatives of
$\B(a,b)=\Gamma(a)\Gamma(b)/\Gamma(a+b)$.  Differentiate those finite local Taylor identities once more and apply the product rule.  The Gamma packet supplies derivative-remainder certificates, so no qualitative theorem about differentiating quotients is required.
\end{proof}

More generally, derivatives of $\log\B$ are finite signed sums of polygamma values:
\[
 \partial_a^r\log\B(a,b)
 =\psi^{(r-1)}(a)-\psi^{(r-1)}(a+b)
 \quad(r\ge1),
\]
and every mixed derivative involving both $a$ and $b$ differentiates only the term $-\log\Gamma(a+b)$.  Derivatives of $\B$ itself are then complete Bell polynomials in these logarithmic derivatives.  For each externally fixed order this is finite polynomial arithmetic and therefore an iPRA theorem schema.

\subsection{Convexity and covariance consequences}

For positive real $a,b$, normalize the beta density
\[
 d\mu(t)=\frac{t^{a-1}(1-t)^{b-1}}{\B(a,b)}\,dt.
\]
This notation abbreviates finite weighted sums and their limit certificates; no probability space is postulated.  Then
\begin{align*}
 \partial_a\log\B(a,b)&=\mathbb E_\mu[\log t],\\
 \partial_b\log\B(a,b)&=\mathbb E_\mu[\log(1-t)],\\
 \partial_a^2\log\B(a,b)&=
 \mathbb E_\mu[(\log t)^2]
 -\mathbb E_\mu[\log t]^2\ge0,\\
 \partial_b^2\log\B(a,b)&\ge0,\\
 \partial_a\partial_b\log\B(a,b)&=
 \mathbb E_\mu[\log t\log(1-t)]
 -\mathbb E_\mu[\log t]\mathbb E_\mu[\log(1-t)].
\end{align*}
The variance inequality is the finite identity
\[
 \left(\sum w_i\right)\left(\sum w_i x_i^2\right)
 -\left(\sum w_i x_i\right)^2
 =\frac12\sum_{i,j}w_iw_j(x_i-x_j)^2\ge0
\]
followed by the integration error bounds.  Thus separate log-convexity of Beta is fully constructive.

\subsection{Derivatives of the multivariate beta function}

For $\boldsymbol\alpha=(\alpha_1,\dots,\alpha_n)$ and
$S=\sum_j\alpha_j$,
\[
 \B(\boldsymbol\alpha)
 =\frac{\prod_j\Gamma(\alpha_j)}{\Gamma(S)}.
\]
For every fixed coordinate $m$,
\begin{equation}
 \boxed{
 \frac{\partial}{\partial\alpha_m}\B(\boldsymbol\alpha)
 =\B(\boldsymbol\alpha)
 \bigl(\psi(\alpha_m)-\psi(S)\bigr).}
 \label{eq:multibeta-derivative}
\end{equation}
This is the source-page formula.  For a variable-length encoded list, the product and sum are primitive-recursive recursors; the coordinate lookup is bounded search.  The derivative order and list length can be parameters of the metatheory, while each internal theorem uses a fixed primitive-recursive implementation.

\section{Multiplication identities and Gosper products}
\label{sec:multiplication}

All identities in this section are finite consequences of Gauss's multiplication theorem
\begin{equation}
 \prod_{j=0}^{n-1}\Gamma\!\left(z+\frac jn\right)
 =(2\pi)^{(n-1)/2}n^{1/2-nz}\Gamma(nz).
 \label{eq:gauss-mult}
\end{equation}
ProveIt's Gamma packet derives \eqref{eq:gauss-mult} from finite Euler approximants, reflection, and a finite sine product, with explicit pole margins.  Once that contract is available, the beta identities require only product regrouping and Gamma recurrence, hence quantifier-free induction.

\subsection{MathWorld's $n$-fold identity}

\begin{theorem}[Beta multiplication]\label{thm:beta-multiplication}
For $n\ge1$ and certified parameters,
\begin{equation}
 \boxed{
 \B(np,nq)=n^{-nq}
 \frac{\displaystyle\prod_{j=0}^{n-1}
       \B\!\left(p+\frac jn,q\right)}
      {\displaystyle\prod_{k=1}^{n-1}\B(kq,q)}.}
 \label{eq:beta-multiplication}
\end{equation}
The denominator product is empty and equals $1$ when $n=1$.
\end{theorem}

\begin{proof}
Apply the beta--gamma identity to the numerator:
\[
 \prod_{j=0}^{n-1}\B\!\left(p+\frac jn,q\right)
 =\Gamma(q)^n
 \frac{\prod_j\Gamma(p+j/n)}
      {\prod_j\Gamma(p+q+j/n)}.
\]
Gauss multiplication makes the quotient
\[
 n^{nq}\frac{\Gamma(np)}{\Gamma(n(p+q))}.
\]
The denominator in \eqref{eq:beta-multiplication} telescopes:
\[
 \prod_{k=1}^{n-1}\B(kq,q)
 =\Gamma(q)^{n-1}
   \prod_{k=1}^{n-1}\frac{\Gamma(kq)}{\Gamma((k+1)q)}
 =\frac{\Gamma(q)^n}{\Gamma(nq)}.
\]
Dividing the first expression by the second gives
$n^{nq}\B(np,nq)$.  For a division-free internal proof, multiply by every displayed Gamma denominator and use the same finite cancellations.
\end{proof}

\subsection{A general product with simultaneous shifts}

Put
\[
 P_N(a,b)=\prod_{i=0}^{N-1}
 \B\!\left(a+\frac iN,b+\frac iN\right),
 \qquad s=a+b.
\]
The two numerator Gamma products in $P_N$ are immediately evaluated by \eqref{eq:gauss-mult}.  The denominator
\[
 D_N(s)=\prod_{i=0}^{N-1}\Gamma\!\left(s+\frac{2i}{N}\right)
\]
behaves differently according to the parity of $N$.

\subsection{Odd modulus}

Let $N=2n+1$.  Multiplication by $2$ permutes the residues modulo $N$.  For $0\le i\le n$ there is no carry in $2i/N$; for $n+1\le i\le2n$ there is one carry, and the residual numerators are $1,3,\dots,2n-1$.  Therefore
\[
 D_N(s)=
 \left[\prod_{r=0}^{N-1}\Gamma\!\left(s+\frac rN\right)\right]
 \prod_{j=0}^{n-1}\left(s+\frac{2j+1}{N}\right).
\]
The first product is Gauss's formula.  The second is
\[
 N^{-n}\prod_{j=0}^{n-1}(Ns+2j+1)
 =N^{-n}2^n\left(\frac{Ns+1}{2}\right)_n.
\]
After cancellation we obtain the first Gosper identity.

\begin{theorem}[Gosper product, odd case]\label{thm:gosper-odd}
For $n\ge1$, $N=2n+1$, and certified parameters,
\begin{equation}
 \boxed{\begin{aligned}
 \prod_{i=0}^{2n}
 \B\!\left(a+\frac{i}{2n+1},b+\frac{i}{2n+1}\right)
 &=\frac{(2n+1)^{(2n+1)/2}\pi^n}{(n-1)!}\\[-1mm]
 &\quad{}\times
 \B\!\left(n,\frac{(2n+1)(a+b)+1}{2}\right)
 \B((2n+1)a,(2n+1)b).
 \end{aligned}}
 \label{eq:gosper-odd}
\end{equation}
\end{theorem}

\begin{proof}
The calculation above gives
\[
 P_N(a,b)=
 \frac{N^{N/2}(2\pi)^n}{\prod_{j=0}^{n-1}(Ns+2j+1)}
 \B(Na,Nb).
\]
If $x=(Ns+1)/2$, the denominator product is $2^n(x)_n$, so
\[
 \frac{(2\pi)^n}{2^n(x)_n}
 =\frac{\pi^n}{(n-1)!}\B(n,x).
\]
Substitution proves \eqref{eq:gosper-odd}.
\end{proof}

For $n=1$, $N=3$, $\B(1,x)=1/x$, and the formula reduces to
\[
 \boxed{
 \B(a,b)\B(a+\tfrac13,b+\tfrac13)
 \B(a+\tfrac23,b+\tfrac23)
 =
 \frac{6\pi\sqrt3}{1+3(a+b)}\B(3a,3b).}
\]

\subsection{Even modulus}

Let $N=2n$.  Now $2i/N=i/n$, so each residue class modulo $n$ occurs once at height $0$ and once at height $1$:
\[
 D_{2n}(s)
 =\prod_{i=0}^{n-1}
   \Gamma\!\left(s+\frac in\right)
   \Gamma\!\left(s+1+\frac in\right).
\]
Use $\Gamma(z+1)=z\Gamma(z)$:
\[
 D_{2n}(s)
 =\left[\prod_{i=0}^{n-1}
   \Gamma\!\left(s+\frac in\right)\right]^2
   \prod_{i=0}^{n-1}\left(s+\frac in\right).
\]
Gauss multiplication and
\[
 \prod_{i=0}^{n-1}\left(s+\frac in\right)
 =n^{-n}\frac{\Gamma(n(s+1))}{\Gamma(ns)}
\]
give
\[
 D_{2n}(s)
 =(2\pi)^{n-1}n^{1-n-2ns}
 \Gamma(ns)\Gamma(n(s+1)).
\]

\begin{theorem}[Gosper product, even case]\label{thm:gosper-even}
For $n\ge1$ and certified parameters,
\begin{equation}
 \boxed{
 \prod_{i=0}^{2n-1}
 \B\!\left(a+\frac{i}{2n},b+\frac{i}{2n}\right)
 =
 \frac{n^n\pi^n\B(n,2n(a+b))\B(2na,2nb)}
 {2^{2n(a+b)-n-1}(n-1)!\,
  \B(n(a+b),n(a+b+1))}.}
 \label{eq:gosper-even}
\end{equation}
\end{theorem}

\begin{proof}
The numerator Gamma products give
\[
 (2\pi)^{2n-1}(2n)^{1-2ns}
 \Gamma(2na)\Gamma(2nb).
\]
Divide by the displayed expression for $D_{2n}(s)$ and simplify the powers:
\[
 P_{2n}(a,b)=
 2^{n+1-2ns}n^n\pi^n
 \frac{\Gamma(2na)\Gamma(2nb)}
 {\Gamma(ns)\Gamma(n(s+1))}.
\]
Expanding the three beta factors on the right side of
\eqref{eq:gosper-even} cancels
$\Gamma(2ns)$ and $\Gamma(n(2s+1))$ and gives exactly this quotient.
\end{proof}

At $n=2$ this is MathWorld's four-factor special case:
\[
 \prod_{i=0}^{3}\B(a+i/4,b+i/4)
 =
 \frac{2^{3-4(a+b)}\pi^2\B(4a,4b)}
 {(a+b)(1+4(a+b))
   \B(2(a+b),2(a+b+1))}.
\]
The equality follows from \eqref{eq:gosper-even} plus the elementary formulas
$\B(2,4s)=1/[4s(4s+1)]$ and Gamma recurrence.

\subsection{Internal proof shape}

The parity split is ideal for iPRA.  Encode the map
\[
 i\longmapsto (2i\bmod N,\ \lfloor2i/N\rfloor)
\]
using the checked division-with-remainder infrastructure.  In the odd case, prove by bounded arithmetic that the residue map is a permutation and that the carry set is $n+1,\dots,2n$.  In the even case, pair $i$ with $i+n$.  Prove product invariance under the odd permutation using ProveIt's checked \code{Core/Permutation} theorem.  All remaining steps are identities of finite products and powers.  The only analytic input is the already isolated Gamma multiplication certificate.

\section{Incomplete and regularized beta functions}
\label{sec:incomplete}

\subsection{Effective definition}

Let $x$ be a real name with $0\le x\le1$.  For $\Rea a>0$, define
\[
 \B_x(a,b)=\int_0^x t^{a-1}(1-t)^{b-1}\,dt.
\]
If $x\le1-\delta$ with $\delta>0$, only the lower endpoint is improper and the compact derivative bound uses the margin $1-t\ge\delta$.  At $x=1$ use the complete-beta construction.  If the sign or endpoint status of a general name $x$ is not decidable, the theorem is stated under one of these positive certificates; no exact-real case split is made.

The regularized incomplete beta function is
\[
 \I_x(a,b)=\frac{\B_x(a,b)}{\B(a,b)}
\]
for positive real parameters, where \cref{thm:basic-laws} supplies a positive lower bound on the denominator over a certified compact domain.

\subsection{Basic identities}

\begin{theorem}[Endpoints, complement, and elementary cases]\label{thm:incomplete-basic}
For positive real parameters and $0\le x\le1$,
\begin{align*}
 \B_0(a,b)&=0,& \B_1(a,b)&=\B(a,b),\\
 \B_x(a,b)+\B_{1-x}(b,a)&=\B(a,b),\\
 \B_x(a,1)&=\frac{x^a}{a},&
 \B_x(1,b)&=\frac{1-(1-x)^b}{b}.
\end{align*}
Consequently
\begin{align*}
 \I_0(a,b)&=0,& \I_1(a,b)&=1,\\
 \I_x(a,b)&=1-\I_{1-x}(b,a),\\
 \I_x(a,1)&=x^a,&
 \I_x(1,b)&=1-(1-x)^b.
\end{align*}
\end{theorem}

\begin{proof}
The endpoint statements are definitions and the tail estimates.  Split the complete integral at $x$ and in the second part substitute $t=1-u$.  The elementary cases use the antiderivatives $t^a/a$ and $-(1-t)^b/b$.  All equalities are first established on rational approximations to $x$ and then transferred using the Lipschitz/error budget.
\end{proof}

\subsection{Recurrences and differentiation in the endpoint}

Integrating the derivative of $t^a(1-t)^b$ on $[0,x]$ gives
\[
 a\B_x(a,b+1)-b\B_x(a+1,b)=x^a(1-x)^b.
\]
Together with the incomplete Pascal identity
\[
 \B_x(a,b)=\B_x(a+1,b)+\B_x(a,b+1),
\]
this yields
\begin{align}
 \B_x(a+1,b)
 &=\frac{a}{a+b}\B_x(a,b)
   -\frac{x^a(1-x)^b}{a+b},
 \label{eq:inc-left-rec}\\
 \B_x(a,b+1)
 &=\frac{b}{a+b}\B_x(a,b)
   +\frac{x^a(1-x)^b}{a+b}.
 \label{eq:inc-right-rec}
\end{align}
After division by the corresponding complete-beta recurrence,
\begin{align}
 \I_x(a+1,b)
 &=\I_x(a,b)-\frac{x^a(1-x)^b}{a\B(a,b)},\\
 \I_x(a,b+1)
 &=\I_x(a,b)+\frac{x^a(1-x)^b}{b\B(a,b)}.
\end{align}

The finite fundamental theorem gives
\begin{align}
 \frac{d}{dx}\B_x(a,b)&=x^{a-1}(1-x)^{b-1},\\
 \frac{d}{dx}\I_x(a,b)&=
 \frac{x^{a-1}(1-x)^{b-1}}{\B(a,b)}.
\end{align}
These derivative statements require a local margin $0<d\le x\le1-d$ unless one uses one-sided endpoint certificates.

Integration by parts in $x$, again interpreted by finite telescoping, gives the source formulas
\begin{align*}
 \int\B_x(a,b)\,dx
 &=x\B_x(a,b)-\B_x(a+1,b)\\
 &=\left(x-\frac{a}{a+b}\right)\B_x(a,b)
   +\frac{x^a(1-x)^b}{a+b},\\
 \int\I_x(a,b)\,dx
 &=x\I_x(a,b)-\frac{a}{a+b}\I_x(a+1,b)\\
 &=\left(x-\frac{a}{a+b}\right)\I_x(a,b)
   +\frac{x^a(1-x)^b}{(a+b)\B(a,b)}.
\end{align*}
An indefinite integral here means the displayed antiderivative identity with a quadratic remainder, not a quotient by constants.

\subsection{Alternative integrals}

For $0<x<1$ with endpoint margins,
\begin{align}
 \B_x(a,b)
 &=2\int_0^{\arcsin\sqrt{x}}
   \sin^{2a-1}\theta\cos^{2b-1}\theta\,d\theta,
 \label{eq:inc-trig}\\
 &=\int_{(1-x)/x}^{\infty}
   \frac{s^{b-1}}{(1+s)^{a+b}}\,ds.
 \label{eq:inc-infinity}
\end{align}
The first uses $t=\sin^2\theta$ and a bisection name for the inverse on the monotone first quadrant.  The second uses $t=1/(1+s)$.  Its upper tail is bounded by $T^{-\Rea a}/\Rea a$.

The special identity
\begin{equation}
 \boxed{\B_x(a,-a)=\frac1a\left(\frac{x}{1-x}\right)^a}
 \label{eq:inc-special}
\end{equation}
follows by differentiating the right side and checking the value at $x=0$.  It is constructive for positive real $x<1$ and a specified positive-real power convention.  For complex parameters and transformed negative arguments, branch data must be included.

\subsection{Polynomial values at positive integers}

\begin{theorem}[Exact polynomial formula]\label{thm:inc-polynomial}
For positive integers $m,n$,
\begin{equation}
 \boxed{
 \B_x(m,n)=
 \sum_{j=0}^{n-1}(-1)^j\binom{n-1}{j}
 \frac{x^{m+j}}{m+j}.}
 \label{eq:inc-poly}
\end{equation}
It is a rational polynomial of degree $m+n-1$.
\end{theorem}

\begin{proof}
Use the finite binomial theorem
$(1-t)^{n-1}=\sum_{j=0}^{n-1}(-1)^j\binom{n-1}{j}t^j$
and integrate each monomial exactly.  This is a finite sum identity with no analytic remainder.
\end{proof}

The complement relation gives a second polynomial centered at $1$.  Multiplying by $(m+n-1)!/[(m-1)!(n-1)!]$ gives exact integer-coefficient formulas for the regularized function after common denominators are cleared.

\subsection{Hypergeometric representation}

For $0\le x\le1-\delta$,
\begin{align}
 \B_x(a,b)
 &=\sum_{n=0}^{\infty}
   \frac{(1-b)_n}{(a+n)n!}x^{a+n}
 \label{eq:inc-series}\\
 &=\frac{x^a}{a}\,
 {}_2F_1(a,1-b;a+1;x).
 \label{eq:inc-hyper}
\end{align}
The equality of the two displayed series uses
$(a)_n/(a+1)_n=a/(a+n)$, an exact finite-product cancellation.  The proof is the compact binomial-series proof of \cref{thm:beta-series}; the geometric margin $x\le1-\delta$ gives an especially simple cutoff.  At $x=1$ one instead uses the endpoint tail and the complete beta value.

A Lean development need not introduce a general hypergeometric-function object.  It may define the right side of \eqref{eq:inc-hyper} as this particular finite-sum name and prove the identity directly.

\subsection{The binomial cumulative-distribution identity}

For $0\le k<n$ and $0\le p\le1$, define the finite polynomial
\[
 F_{n,k}(p)=\sum_{j=0}^{k}\binom nj p^j(1-p)^{n-j}.
\]
No probability theory is needed.

\begin{theorem}[Binomial CDF identity]\label{thm:binomial-cdf}
\[
 \boxed{
 F_{n,k}(p)=\I_{1-p}(n-k,k+1)
 =1-\I_p(k+1,n-k).}
\]
\end{theorem}

\begin{proof}
The two regularized beta expressions are equal by complementarity.  Differentiate the finite polynomial.  Using
$j\binom nj=n\binom{n-1}{j-1}$ and
$(n-j)\binom nj=n\binom{n-1}{j}$, adjacent terms telescope:
\[
 F'_{n,k}(p)
 =-\frac{n!}{k!(n-k-1)!}p^k(1-p)^{n-k-1}.
\]
By the factorial formula for $\B(k+1,n-k)$, this is exactly the derivative of
$1-\I_p(k+1,n-k)$.  Both sides equal $1$ at $p=0$.  The finite fundamental theorem proves equality.  Alternatively, induction on $n$ using Pascal's identity stays entirely in polynomial arithmetic.  The omitted boundary case $k=n$ has $F_{n,n}(p)=1$ by the binomial theorem; it is handled separately because the regularized parameter $n-k$ would be zero.
\end{proof}

\subsection{Continued fraction: exact finite target and remaining lemma}

The source gives
\[
 \B_x(a,b)=
 \frac{x^a(1-x)^b}
 {a\left(1+\cfrac{d_1}{1+\cfrac{d_2}{1+\cdots}}\right)}
\]
with
\[
 d_{2m+1}=
 -\frac{(a+m)(a+b+m)x}{(a+2m)(a+2m+1)},
 \qquad
 d_{2m}=
 \frac{m(b-m)x}{(a+2m-1)(a+2m)}.
\]
Define continuants by
\[
 P_{-1}=1,\ P_0=1,\quad
 Q_{-1}=0,\ Q_0=1,\quad
 P_n=P_{n-1}+d_nP_{n-2},\quad
 Q_n=Q_{n-1}+d_nQ_{n-2}.
\]
The $n$th finite continued fraction is $P_n/Q_n$, and its determinant identity
\[
 P_nQ_{n-1}-P_{n-1}Q_n
 =(-1)^{n-1}d_1d_2\cdots d_n
\]
is a one-line induction, fully internal to iPRA.

The remaining analytic theorem is a positive-remainder representation comparing
\[
 \frac{x^a(1-x)^bQ_n}{aP_n}
\]
with $\B_x(a,b)$.  Classical proofs obtain it from Gauss contiguous relations for ${}_2F_1$; the signs of the remainder yield the stated $4m/4m+1$ lower and $4m+2/4m+3$ upper convergents.  A complete iPRA proof should:
\begin{enumerate}
\item express every contiguous relation as a finite coefficient identity for truncated hypergeometric sums;
\item attach the geometric tail from $x\le1-\delta$;
\item derive the remainder as an integral with a nonnegative rational-mesh integrand for $a,b>0$;
\item use the complement relation when $x$ is near $1$.
\end{enumerate}
The continuant algebra and formal target are complete; the general alternating-remainder lemma is classified \statusC{} in this packet.

\section{Multivariate beta and fixed-dimensional simplex integrals}
\label{sec:multivariate}

For a finite nonempty list
$\boldsymbol\alpha=(\alpha_1,\ldots,\alpha_n)$ with positive real parts, define
\[
 \B(\boldsymbol\alpha)
 =\frac{\prod_{j=1}^{n}\Gamma(\alpha_j)}
        {\Gamma(S)},
 \qquad S=\sum_{j=1}^{n}\alpha_j.
\]
For $n=2$ this is the ordinary beta function.

\subsection{Coordinate recurrences and Pascal's identity}

\begin{theorem}[Multivariate recurrence]\label{thm:multi-recurrence}
For each $j$,
\[
 \boxed{
 \B(\alpha_1,\ldots,\alpha_j+1,\ldots,\alpha_n)
 =\frac{\alpha_j}{S}\B(\boldsymbol\alpha).}
\]
Consequently
\[
 \boxed{
 \B(\boldsymbol\alpha)=
 \sum_{j=1}^{n}
 \B(\alpha_1,\ldots,\alpha_j+1,\ldots,\alpha_n).}
\]
\end{theorem}

\begin{proof}
Apply $\Gamma(z+1)=z\Gamma(z)$ in the numerator and denominator.  Summing the coordinate formulas multiplies $\B(\boldsymbol\alpha)$ by
$\sum_j\alpha_j/S=1$.  For internalization, clear the common denominator $S$ before the finite sum.  The list sum and product are primitive-recursive recursors.
\end{proof}

Repeated shifts give, for a vector $\mathbf m\in\N^n$,
\[
 \B(\boldsymbol\alpha+\mathbf m)
 =
 \frac{\prod_j(\alpha_j)_{m_j}}
      {(S)_{\sum_jm_j}}
 \B(\boldsymbol\alpha).
\]
At positive integer arguments this becomes
\[
 \B(m_1,\ldots,m_n)
 =\frac{\prod_j(m_j-1)!}
        {(\sum_jm_j-1)!}.
\]
This is the direct analogue of the multinomial coefficient relation.

\subsection{Derivative formula}

The source-page derivative
\[
 \frac{\partial}{\partial\alpha_j}\B(\boldsymbol\alpha)
 =
 \B(\boldsymbol\alpha)\bigl(\psi(\alpha_j)-\psi(S)\bigr)
\]
was proved in \eqref{eq:multibeta-derivative}.  Higher derivatives are again finite Bell polynomials in polygamma differences.  For each fixed order and list length, this is a finite formal identity.  A variable list can be represented by a code and a length; the internal theorem quantifies only over the code and uses bounded lookup.

\subsection{Simplex integral}

For fixed $n\ge2$, let
\[
 \Delta_{n-1}=
 \{(t_1,\ldots,t_{n-1})\in[0,1]^{n-1}:
   t_1+\cdots+t_{n-1}\le1\},
 \qquad
 t_n=1-\sum_{j<n}t_j.
\]
Then
\begin{equation}
 \boxed{
 \int_{\Delta_{n-1}}\prod_{j=1}^{n}t_j^{\alpha_j-1}
 \,dt_1\cdots dt_{n-1}
 =\B(\boldsymbol\alpha).}
 \label{eq:simplex-beta}
\end{equation}
A constructive proof avoids a general multivariate integration theorem.  Integrate $t_{n-1}$ first and use the substitution
\[
 t_{n-1}=(1-t_1-\cdots-t_{n-2})u.
\]
The inner integral becomes
\[
 (1-t_1-\cdots-t_{n-2})^{\alpha_{n-1}+\alpha_n-1}
 \B(\alpha_{n-1},\alpha_n).
\]
Repeat.  After $n-2$ steps, the product of beta factors telescopes to the Gamma quotient.  Every stage is a one-dimensional finite Riemann sum with endpoint bounds; the dimension $n$ is externally fixed, or encoded by a simultaneous primitive recursion if a uniform family is desired.

\subsection{Dirichlet normalization without probability axioms}

For positive real parameters, dividing the simplex integrand by
$\B(\boldsymbol\alpha)$ produces a finite-sum name whose integral is $1$.  Expectations of monomials are merely ratios of shifted multivariate beta values:
\[
 \frac{1}{\B(\boldsymbol\alpha)}
 \int_{\Delta_{n-1}}
 \prod_j t_j^{\alpha_j+m_j-1}\,dt
 =
 \frac{\prod_j(\alpha_j)_{m_j}}
      {(S)_{\sum_jm_j}}.
\]
Thus the algebraic moment formulas usually stated for the Dirichlet distribution are already iPRA theorems about finite weighted sums and their certified limits; no measure-theoretic probability space is required.

\section{Quantitative asymptotics}
\label{sec:asymptotics}

An asymptotic statement is meaningful in iPRA only after it is replaced by a theorem with explicit integer thresholds and rational error bounds.  The Gamma packet supplies effective Stirling formulas with signed remainders.  We record the beta consequences in a form that is directly formalizable.

\subsection{Both parameters large}

Suppose $x,y\ge1$ are positive real names.  Use the effective logarithmic Stirling representation
\[
 \Gamma(u)=\sqrt{2\pi}\,
 u^{u-1/2}e^{-u}e^{\rho(u)},
 \qquad
 0<\rho(u)<\frac{1}{12u}
 \quad(u\ge1).
\]
Substitution in the beta--gamma quotient gives the exact form
\begin{equation}
 \B(x,y)=
 \sqrt{2\pi}\,
 \frac{x^{x-1/2}y^{y-1/2}}
      {(x+y)^{x+y-1/2}}
 e^{E(x,y)},
 \qquad
 E=\rho(x)+\rho(y)-\rho(x+y).
 \label{eq:beta-stirling-exact}
\end{equation}
Hence
\begin{equation}
 |E(x,y)|<
 \frac1{12x}+\frac1{12y}+\frac1{12(x+y)}
 =:\eta(x,y).
 \label{eq:beta-stirling-error}
\end{equation}
Since $|e^E-1|\le2|E|$ when $\eta\le1$, we obtain a fully quantitative version of the Wikipedia asymptotic:
\begin{equation}
 \left|
 \frac{\B(x,y)}
 {\sqrt{2\pi}\,
 x^{x-1/2}y^{y-1/2}/(x+y)^{x+y-1/2}}
 -1\right|
 \le2\eta(x,y).
 \label{eq:beta-stirling-relative}
\end{equation}
For a requested relative error $2^{-p}$ it suffices, for example, to require
$x,y\ge2^{p+5}$.  The bound is deliberately coarse.  All powers are positive-real powers and all logarithms are evaluated by the elementary primitive-recursive algorithms.

\subsection{A fixed ratio}

Let $x=\lambda N$, $y=(1-\lambda)N$ with a supplied rational margin
$d\le\lambda\le1-d$.  Rearranging \eqref{eq:beta-stirling-exact} gives
\[
 \B(\lambda N,(1-\lambda)N)
 =
 \sqrt{\frac{2\pi}{N}}\,
 \lambda^{\lambda N-1/2}
 (1-\lambda)^{(1-\lambda)N-1/2}
 e^{E}.
\]
The error satisfies $|E|\le C_d/N$ with a rational $C_d$.  This form isolates the exponential entropy factor and is often more useful in combinatorics.  It is a theorem schema for each fixed named $\lambda$ and its margins.

\subsection{One parameter fixed}

The source states
\[
 \B(x,y)\sim\Gamma(y)x^{-y}
 \qquad(x\to\infty,\ y\text{ fixed}).
\]
A particularly clean iPRA version is on the natural grid.

\begin{theorem}[Fixed-parameter beta limit]\label{thm:fixed-y-asymp}
Fix a named $y$ with $\Rea y\ge\alpha>0$ and $|y|\le d$.  Then
\[
 N^y\B(y,N+1)\longrightarrow\Gamma(y)
\]
with the explicit modulus in \eqref{eq:beta-gamma-limit} (take $b=1$).  By symmetry,
\[
 \boxed{\B(N+1,y)=\Gamma(y)N^{-y}(1+\varepsilon_N),}
\]
where $|\varepsilon_N|$ is bounded by the right side of
\eqref{eq:beta-gamma-limit} divided by a supplied positive lower bound for $\Gamma(y)$.
\end{theorem}

For a general real $x\to\infty$, supply an integer $N$ and a rational bracket
$N\le x\le N+1$.  On positive real inputs, differentiation under the integral gives
\[
 \partial_x\log\B(x,y)
 =\frac{\int_0^1 t^{x-1}(1-t)^{y-1}\log t\,dt}
        {\B(x,y)}<0,
\]
so Beta is decreasing in $x$.  Bracketing between the two neighboring natural-grid values and using
$(1+1/N)^y=1+O(1/N)$ transfers the limit effectively.  The comparison can alternatively be obtained directly from the Gamma ratio and Stirling.

\subsection{Higher-order expansions}

For every fixed $M$, the Gamma packet's finite Euler--Maclaurin theorem gives
\[
 \log\Gamma(z)
 =
 \left(z-\frac12\right)\log z-z+\frac12\log(2\pi)
 +\sum_{r=1}^{M-1}
   \frac{B_{2r}}{2r(2r-1)z^{2r-1}}
 +R_M(z),
\]
with a signed rational majorant for $R_M$ on a positive real ray.  Taking the combination
\[
 \log\B(x,y)=\log\Gamma(x)+\log\Gamma(y)-\log\Gamma(x+y)
\]
produces an arbitrary fixed-order beta expansion.  Its coefficients are finite rational expressions in Bernoulli numbers and reciprocal powers of $x,y,x+y$.  No new analytic argument is needed: the three remainders are added.  In Lean, $M$ is an external theorem parameter or a natural parameter of a primitive-recursive coefficient generator; the induction invariant is the finite Euler--Maclaurin telescoper already required by the Gamma development.

\subsection{Laplace concentration as a finite estimate}

For large positive $x,y$, the integrand
\[
 \exp\!\bigl((x-1)\log t+(y-1)\log(1-t)\bigr)
\]
has its maximum near $t=x/(x+y)$.  A direct constructive Laplace proof can split the interval into a rational central window and two tails.  Concavity follows from the explicit negative second derivative
\[
 -\frac{x-1}{t^2}-\frac{y-1}{(1-t)^2}.
\]
Finite Taylor inequalities give a Gaussian central bracket and geometric tail decay.  This yields \eqref{eq:beta-stirling-relative} independently of Gamma.  The argument is useful for future work on beta distributions, but the Gamma/Stirling route is shorter and is the one chosen for the formalization dependency graph.

\section{Complex continuation, contour formulas, and reciprocal-beta integrals}
\label{sec:complex}

\subsection{Meromorphic continuation through Gamma}

The right-half-plane integral is only the initial definition.  The quotient
\[
 \B(a,b)=\Gamma(a)\Gamma(b)\,\Gamma(a+b)^{-1}
\]
extends Beta meromorphically to pairs of complex names on any certified compact region avoiding the relevant pole divisors.  Constructively, ``avoiding'' means that the input comes with lower bounds for its distances from the finitely many nonpositive integers in the region.  The reciprocal Gamma function is entire in the effective sense of the Gamma packet, so the cleared definition
\[
 \B(a,b)\Gamma(a+b)=\Gamma(a)\Gamma(b)
\]
remains preferable near cancellations.

The possible singular divisors are $a\in\{0,-1,-2,\dots\}$ and
$b\in\{0,-1,-2,\dots\}$; zeros of $1/\Gamma(a+b)$ can cancel some intersections.  A full local pole-order classification is finite once the relevant integers are bounded by the radius certificate.

\subsection{The Pochhammer contour}

The source states
\begin{equation}
 (1-e^{2\pi i\alpha})(1-e^{2\pi i\beta})\B(\alpha,\beta)
 =
 \int_C t^{\alpha-1}(1-t)^{\beta-1}\,dt,
 \label{eq:pochhammer}
\end{equation}
where $C$ is a Pochhammer contour around $0$ and $1$.  The formula gives analytic continuation because the contour stays away from the branch points.

A constructive implementation can be made as follows.

\begin{enumerate}
\item Replace $C$ by a rational polygon consisting of two small circles approximated by regular polygons, joined by upper and lower rational-offset segments.
\item On each segment carry explicit branches of $\log t$ and $\log(1-t)$, together with rational bounds on their derivatives.
\item Define the contour integral by finite line-segment Riemann sums.
\item Pair the four traversals of the connecting interval.  Their branch multipliers are
$1$, $e^{2\pi i\alpha}$, $e^{2\pi i\beta}$, and
$e^{2\pi i(\alpha+\beta)}$; finite algebra leaves the prefactor in
\eqref{eq:pochhammer}.
\item Bound the two small loops by
$C\rho^{\Rea\alpha}$ and $C\rho^{\Rea\beta}$ in the initial right-half-plane domain.  This proves \eqref{eq:pochhammer} there.
\item Both sides have effective local power-series certificates off the integer divisors.  Continue along a finite chain of rational disks, using coefficient uniqueness on overlaps.
\end{enumerate}

Each step has a clear primitive-recursive realization, but the archive does not instantiate the polygonal branch and disk-chain certificates.  The Pochhammer formula is therefore \statusC{} here.  Meromorphic continuation through the Gamma quotient is already \statusB{} and suffices for all algebraic identities in the report.

\subsection{Reciprocal-beta trigonometric integrals}

The source displays, among others,
\begin{equation}
 \int_0^{\pi/2}\cos^{x-1}\theta\cos(y\theta)\,d\theta
 =
 \frac{\pi}
 {2^x x\,
 \B\!\left(\frac{x+y+1}{2},\frac{x-y+1}{2}\right)}.
 \label{eq:reciprocal-beta-cos}
\end{equation}
Similar formulas over $[0,\pi]$ include factors
$\sin(\pi y/2)$ or $\cos(\pi y/2)$.  The ordinary finite-interval integral in
\eqref{eq:reciprocal-beta-cos} needs only $\Rea x>0$: for a fixed complex
$y$, the cosine factor is bounded on $[0,\pi/2]$.  The stronger strip
$\Rea x>|\Rea y|-1$ puts both displayed Beta arguments in the right half-plane
and therefore supports a direct Beta-integral derivation.  Outside that strip,
the reciprocal Beta on the right is read through its entire reciprocal-Gamma
factors; no pole division is needed.

For an integer frequency $y=m$, \eqref{eq:reciprocal-beta-cos} has a fully finite proof.  The multiple-angle recurrence expresses $\cos(m\theta)$ as a polynomial in $\cos\theta$; after multiplying by $\cos^{x-1}\theta$, every term is a trigonometric beta integral.  Gamma recurrence and duplication simplify the resulting finite sum to the right side.  The induction on $m$ and the polynomial coefficient identity are quantifier-free.

For arbitrary named real or complex $y$, one route is:
\[
 \cos(y\theta)=
 {}_2F_1\!\left(-\frac y2,\frac y2;\frac12;\sin^2\theta\right)
\]
or an equivalent exponential representation, followed by the beta integral and Gauss summation at $1$.  A constructive proof requires a uniform hypergeometric tail and a pole-margin version of Gauss summation.  Those contracts are not yet present in the inspected ProveIt arithmetic library, so the general-frequency formulas are \statusC{}; the integer-frequency family is \statusA/\statusB{}.

\subsection{Reflection identities outside the integral strip}

The formula
\[
 \B(x,1-x)=\frac{\pi}{\sin\pi x}
\]
is an ordinary integral identity only for $0<\Rea x<1$.  Outside that strip it is a meromorphic identity and requires a distance certificate from the integers before division by sine.  A bare hypothesis $x\notin\Z$ is not constructive enough.  Given a radius $R$ and
\[
 |x-k|\ge\delta>0
 \quad\text{for }|k|\le2R+1,
\]
the effective sine product supplies a rational lower bound for
$|\sin\pi x|$, and guarded inversion completes the formula.

\subsection{Branches in incomplete-beta transformations}

The source also writes
\[
 \B_x(a,b)=(-1)^a
 \B_{x/(x-1)}(a,1-a-b).
\]
For noninteger complex $a$, the notation $(-1)^a$ is branch-dependent, and the transformed upper endpoint is negative for $0<x<1$.  A correct constructive statement must specify a branch of $\log(-1)$ and a path avoiding the branch points.  For integer $a$, the sign is unambiguous and the formula follows by the rational substitution $u=t/(t-1)$ on finite intervals.  The unqualified complex formula is classified \statusD{} as written, while a branch-certified version is \statusC{}.

\subsection{Cleared identities as the preferred complex interface}

When analytic continuation introduces poles, the most robust internal targets are cleared equations:
\begin{align*}
 \Gamma(a+b)\B(a,b)&=\Gamma(a)\Gamma(b),\\
 \sin(\pi x)\B(x,1-x)&=\pi,\\
 \pi\binom nk\prod_{j=0}^n(k-j)
 &=(-1)^n n!\sin(\pi k).
\end{align*}
These formulas do not perform undefined division at removable singularities.  Local meromorphic statements are recovered only after the caller supplies the relevant apartness margins.

\section{Source audit: qualifications, corrections, and non-theorem material}
\label{sec:source-audit}

The goal is not merely to reproduce formulas but to determine what can be proved with their stated hypotheses.  The following audit is part of the mathematical result.

\subsection{Corrections and required qualifications}

\begin{longtable}{@{}p{0.25\textwidth}p{0.66\textwidth}@{}}
\toprule
Source item & Constructive audit\\
\midrule\endhead
Wikipedia symmetric series for $\B(z,z)$ &
False as rendered with the explicitly stated rising-factorial convention.  At $z=1$ the first series has terms not tending to zero; the second Abel expression equals $2/3$, not $1$.  See \cref{warn:false-series}.\\

Fractional-linear integral with parameter $c\le1$ &
The displayed substitution requires $c<1$.  At $c=1$ it degenerates and the ordinary integrand/prefactor formula is undefined.  A limiting statement would need separate hypotheses.\\

Central radical formula with $z\ne1$ &
For positive real $z$ the proof works also at $z=1$; the exclusion is unnecessary.  For complex $z$, roots require branch choices that the source does not state.\\

Reflection $B(x,1-x)=\pi/\sin\pi x$ for $x\notin\Z$ &
The integral proof requires $0<\Rea x<1$.  Elsewhere it is meromorphic.  Constructive division requires a positive distance from the relevant integers, not merely negated membership.\\

Binomial interpolation quotient &
The raw quotient has $0/0$ at integer $k$.  Use the cleared entire identity or the factorial value at integers.\\

Incomplete-beta transformation with $(-1)^a$ &
Branch-ambiguous for noninteger complex $a$ and a negative transformed endpoint.  It is directly valid for integer $a$, or with an explicit branch/path certificate.\\

Continued-fraction bounds &
The finite continuant algebra is elementary.  The alternating inequalities and rapid-convergence statement require positivity and a remainder theorem; they are not consequences of formal continued-fraction notation alone.\\

MathWorld factorial notation &
Expressions such as $(p-1)!$ are literal factorials only at positive integers.  For general complex parameters they abbreviate Gamma values.\\

MathWorld polar derivation &
Correct classically, but unnecessary for the iPRA target.  A finite polar-grid proof needs a Jacobian error analysis; the one-dimensional beta-to-gamma limit is simpler.\\

Software and application paragraphs &
Statements about Excel, SciPy, string scattering, or preferential attachment are factual/application claims, not beta-function identities.  They are not formalized in this packet.\\
\bottomrule
\end{longtable}

\subsection{The corrected counterexample calculation}

For clarity, at $z=1$ the second displayed symmetric expression is
\[
 2\sum_{n\ge0}\frac{(-2)_n r^n}{(n+1)n!}.
\]
The rising factorial terminates:
\[
 (-2)_0=1,\qquad(-2)_1=-2,\qquad(-2)_2=2,\qquad(-2)_n=0\ (n\ge3).
\]
Therefore the $r\to1^-$ limit is
\[
 2\left(1-1+\frac13\right)=\frac23.
\]
The exact-rational evidence file checks this arithmetic.  This corrects an earlier scratch calculation and is the value used throughout the final packet.

\subsection{Statements not fully proved here}

The following are mathematically plausible standard results but remain at status \statusC{} in this packet:

\begin{enumerate}
\item the Pochhammer-contour identity with explicit branch and polygonal-contour certificates;
\item the reciprocal-beta trigonometric integrals for arbitrary noninteger frequency;
\item the general alternating-convergent remainder theorem for the incomplete-beta continued fraction;
\item a completely internalized variable-dimension simplex integration theorem;
\item a kernel-checked production of all analytic statements as \code{Prf PRAAx} terms.
\end{enumerate}

These are not omitted because of classical logic.  Each has a finitary route; the missing work is the construction and checking of substantial quantitative certificates.

\subsection{Why the audit matters for iPRA}

Classical expositions often use analytic continuation to make a formula true after silently changing its meaning.  That practice is especially dangerous in a weak arithmetic setting.  An integral, a meromorphic quotient, a branch-dependent contour value, and a removable-singularity extension are four different algorithms.  They may agree on overlapping domains, but the agreement itself needs a theorem and a modulus.  This report therefore states domains at every transition and favors cleared equations at possible singularities.

\section{Lean and ConciseIPRA formalization plan}
\label{sec:lean}

\subsection{Repository interfaces}

The intended target is the checked embedded calculus used by ProveIt's
\code{IPRADerivations} library.  A finished theorem has the shape
\[
 ([]:\operatorname{List}(\code{Fm PTm 0}))
 \ \vdash[\code{PRAAx}]\ \varphi.
\]
The checked arithmetic modules already provide:

\begin{itemize}
\item primitive-recursive definitions and defining equations;
\item semiring laws, cancellation, order, and decidable rational-code comparisons;
\item finite sums, finite products, powers, and factorials;
\item pairing, simultaneous recursion, division with remainder, and bounded search;
\item product invariance under an explicitly invertible permutation;
\item a tactic layer for polynomial identities over encoded integers.
\end{itemize}

Once an ordinary iPRA derivation exists, \code{pra_iff_native} in
\code{ConciseIPRA.NativeBiInterpretation} transports it to the native scalar core.  The bridge does not create the derivation; it is the final checked translation.

\subsection{Host-level organization types}

The companion file \code{lean/BetaIPRAOutline.lean} uses ordinary Lean structures to organize obligations.  These structures are metatheoretic bookkeeping, not new sorts added to iPRA.

\begin{lstlisting}[language={}]
structure RatCode where
  num : Int
  den : Nat
  den_pos : 0 < den

structure CQRat where
  re : RatCode
  im : RatCode

structure RegularName where
  approx : Nat -> RatCode
  regular : forall p q,
    absQ (approx p - approx q) <= dyadic p + dyadic q

structure BetaDomain (a b : CName) where
  A B R : Nat
  A_pos : 0 < A
  B_pos : 0 < B
  re_a_margin : Re a >= 1 / A
  re_b_margin : Re b >= 1 / B
  norm_a : norm1 a <= R
  norm_b : norm1 b <= R
\end{lstlisting}

In the actual object theory, a \code{RegularName} is not quantified over.  A particular \code{approx} implementation is compiled to a primitive-recursive symbol, and its regularity theorem is an assumption or previously derived theorem about that symbol.

\subsection{The concrete beta evaluator}

For domain data $D=(A,B,R)$ and precision $p$, define:
\begin{align*}
 L(D,p)&=p+A+B+10,\\
 h(D,p)&=AB\,L(D,p),\\
 \varepsilon(D,p)&=2^{-h(D,p)},\\
 C(D,p)&=2(R+1)\varepsilon(D,p)^{-2R-3},\\
 m(D,p)&=2^{p+7}(1+\lceil C(D,p)\rceil).
\end{align*}
The evaluator is the left sum
\[
 \operatorname{betaApprox}(D,a,b,p)
 =\frac{1-2\varepsilon}{m}
 \sum_{j=0}^{m-1}
 f_{a,b}\!\left(
 \varepsilon+\frac{j(1-2\varepsilon)}{m}\right),
\]
with each elementary-function call made at the precision allocated by the product-error lemma.  All divisions are rational-code divisions with visibly positive denominators.

The first Lean milestones are:

\begin{enumerate}
\item prove that every sample point lies in
$[\varepsilon,1-\varepsilon]$;
\item prove the endpoint bounds
\eqref{eq:lefttail}--\eqref{eq:righttail};
\item prove the derivative majorant $C(D,p)$;
\item prove that the Riemann error is below $2^{-p-7}$;
\item prove regularity and independence from enlarged valid domain data.
\end{enumerate}

The internal formula for regularity is a cleared integer inequality.  It contains no reference to an abstract integral.

\subsection{Module dependency order}

\begin{longtable}{@{}p{0.21\textwidth}p{0.46\textwidth}p{0.25\textwidth}@{}}
\toprule
Planned module & Main declarations & iPRA invariant\\
\midrule\endhead
\code{Beta/RatCode} &
Signed fractions, dyadics, complex pairs, absolute-value comparisons, error propagation &
Integer cross products; positivity of denominators\\

\code{Beta/Elementary} &
Certified \code{exp}, \code{log}, positive powers, sine, cosine, square root and positive rational powers &
Finite Taylor/geometric sums and factorial tails\\

\code{Beta/Names} &
Regular-name equality, arithmetic, guarded inverse, substitution of fixed names &
Precision fixed before induction; rational inequalities only\\

\code{FiniteIntegral} &
Riemann sums, common refinement, finite fundamental theorem, integration by parts, substitution &
Telescoping sums with a quadratic remainder\\

\code{Beta/Complete} &
\code{betaApprox}, endpoint and mesh certificates, symmetry, positivity &
Mesh-index induction and list reversal\\

\code{Beta/Recurrence} &
Pascal, balanced integration by parts, coordinate shifts, factorial values &
Finite product concatenation and endpoint terms\\

\code{Beta/GammaBridge} &
Effective limit \eqref{eq:beta-gamma-limit}, exact Euler ratio, beta--gamma &
Finite logarithm inequalities and Gamma approximant error\\

\code{Beta/Transforms} &
Trigonometric, infinite-interval, power, fractional-linear, central radical forms &
Substitution telescope plus endpoint budgets\\

\code{Beta/SeriesProduct} &
Binomial series, Euler product, product tails &
Finite Taylor polynomial; product telescope; $\sum k^{-2}$ bound\\

\code{Beta/Derivatives} &
Logarithmic moments and digamma/polygamma formulas &
Finite parameter Taylor polynomial\\

\code{Multiplication} &
General multiplication and Gosper odd/even products &
Permutation of residues; carry-count product\\

\code{Beta/Incomplete} &
Incomplete evaluator, recurrences, polynomial values, hypergeometric name, binomial CDF &
Finite interval split; binomial theorem; polynomial telescoping\\

\code{ContinuedFraction} &
Continuants, determinant identity, positive remainder &
Second-order recursion and sign invariant\\

\code{Beta/Multi} &
List Gamma quotient, shifts, Pascal, fixed-dimensional simplex &
List sum/product recursion; iterated one-dimensional beta reduction\\

\code{Beta/Asymptotic} &
Stirling combination, fixed-parameter limit &
Addition of three explicit Gamma remainders\\

\code{Beta/ObjectTheory} &
Reified \code{Fm PTm} statements and \code{Prf PRAAx} terms &
Only \code{PRA.qfInd} on encoded finite invariants\\
\bottomrule
\end{longtable}

\subsection{Representative theorem statements in Lean}

The metatheoretic analysis layer can expose the following shapes.  The exact names depend on the checked-out ProveIt revision.

\begin{lstlisting}[language={}]
theorem beta_symm
    (D : BetaDomain a b) :
    NameEq (betaName D a b) (betaName D.swap b a) := by
  -- reverse the common Riemann mesh; discharge tails
  ...

theorem beta_pascal
    (D : BetaDomain a b) :
    NameEq (betaName D a b)
      (betaName (D.shiftLeft) (a+1) b
       + betaName (D.shiftRight) a (b+1)) := by
  -- pointwise identity t + (1-t) = 1
  ...

theorem beta_shift_left
    (D : BetaDomain a b) :
    NameEq (betaName D.shiftLeft (a+1) b)
      (a / (a+b)) * betaName D a b := by
  -- Pascal plus finite integration by parts
  ...

theorem beta_gamma
    (D : BetaDomain a b)
    (G : EffectiveGammaContract) :
    NameEq (betaName D a b)
      (G.gamma a * G.gamma b / G.gamma (a+b)) := by
  -- effective N^a B(a,b+N) limit and exact Euler ratio
  ...

theorem incomplete_beta_nat
    (m n : Nat) (hm : 0 < m) (hn : 0 < n) :
    PolyEq (incBetaPoly m n)
      (sum n (fun j =>
        (-1)^j * choose (n-1) j * X^(m+j) / (m+j))) := by
  -- finite binomial theorem
  ...
\end{lstlisting}

The delivered outline keeps the analytic holes visible as \code{sorry}; it is not represented as a completed Lean proof.

\subsection{Object-theory formulas}

A name equality must be turned into an arithmetical sentence.  Let
$\mathsf{BA}(D,a,b,p)$ be the rational-code output of the evaluator and let
$\operatorname{close}(u,v,p)$ be the Boolean relation
$|u-v|\le2^{-p}$ after positive denominators are cleared.  The symmetry target is
\[
 \forall D,a,b,p\;
 \operatorname{Dom}(D,a,b)\to
 \operatorname{close}\bigl(
 \mathsf{BA}(D,a,b,p+c),
 \mathsf{BA}(D^\mathrm{op},b,a,p+c),p\bigr).
\]
Here $D,a,b$ are finite tuples of natural-number codes, and
$\operatorname{Dom}$ is a primitive-recursive Boolean conjunction.

The Pascal target is
\[
 \operatorname{close}\bigl(
 \mathsf{BA}(a,b,p+c),
 \mathsf{BA}(a+1,b,p+c)+
 \mathsf{BA}(a,b+1,p+c),p\bigr).
\]
At a common mesh, the summand identity is exact.  The proof object consists of:
\begin{enumerate}
\item the derivation of the summand identity by the ring tactic;
\item finite-sum congruence;
\item the endpoint and mesh inequalities;
\item integer triangle inequalities.
\end{enumerate}

For the beta--gamma identity, division is avoided:
\[
 \operatorname{close}\bigl(
 \mathsf{BA}(a,b,p+c)\cdot\mathsf{GA}(a+b,p+c),
 \mathsf{GA}(a,p+c)\cdot\mathsf{GA}(b,p+c),p\bigr).
\]
Only after this theorem is available does a guarded-inverse lemma derive the quotient form.

\subsection{Quantifier-free induction invariants}

The important invariants are concrete.

\begin{description}
\item[Mesh reversal.]  The sum of the first $r$ beta samples for $(a,b)$ equals the sum of the last $r$ reversed samples for $(b,a)$, up to a tracked rational evaluation error.
\item[Pascal.]  The difference between the accumulated left side and the two accumulated right sides is exactly zero for the first $r$ samples.
\item[Integration by parts.]  The accumulated discrete derivative equals the endpoint difference plus a remainder bounded by $rCh^2$.
\item[Rising products.]  After $r$ shifts,
$(a+b)_r\B(a+r,b)=(a)_r\B(a,b)$ at the chosen finite precision.
\item[Odd Gosper permutation.]  The multiplicity counter for residue $j$ among $2i\bmod(2n+1)$ is $1$ after all indices have been processed; the carry product contains exactly the odd residues.
\item[Incomplete polynomial.]  The derivative of the partial polynomial through term $r$ equals the first $r$ terms of the finite binomial expansion.
\item[Binomial CDF.]  The derivative of the first $r$ summands telescopes to two boundary terms; after $r=k+1$ only the final boundary remains.
\end{description}

Each invariant is an equality or inequality among primitive-recursive natural-number terms.  No induction formula contains the unbounded precision quantifier used in name equality.

\subsection{Use of the checked ProveIt arithmetic library}

The most direct imports for the eventual implementation are expected to include
\begin{lstlisting}[language={}]
import ConciseIPRA.NativeBiInterpretation
import IPRADerivations.Core.Basic
import IPRADerivations.Core.Definitions
import IPRADerivations.Core.Sums
import IPRADerivations.Core.Order
import IPRADerivations.Core.Powers
import IPRADerivations.Core.Division
import IPRADerivations.Core.Search
import IPRADerivations.Core.Permutation
\end{lstlisting}
Exact module paths must be adjusted to the current workspace.  The Gamma packet's host-level \code{GammaIPRA.lean} is useful as an interface sketch, but its analytic certificate structures are uninstantiated and should not be imported as if they were proved theorems.

\subsection{Definition of done}

A complete formalization should satisfy all of the following.

\begin{enumerate}
\item Every theorem exported as a beta identity has an explicit term of type
$[]\vdash[\code{PRAAx}]\varphi$.
\item A \texttt{\#print axioms} audit shows no \code{sorryAx}; standard Lean metatheory axioms are listed separately.
\item Every analytic evaluator is a fixed primitive-recursive symbol with defining equations.
\item Every use of inverse includes a proved positive margin.
\item Every improper integral, infinite sum, product, derivative, or asymptotic statement has a closed cutoff function.
\item The final native theorem is obtained only by applying the checked
\code{pra_iff_native} bridge to the already constructed ordinary derivation.
\end{enumerate}

The current archive meets the mathematical-specification and outline requirements, not this final kernel-checked definition of done.

\section{Formula-by-formula coverage ledger}
\label{sec:inventory}

The table below records the main mathematical statements appearing on the two requested pages.  Closely related display lines are grouped when their proofs are identical.  The CSV artifact contains the same ledger in machine-readable form.

\begin{longtable}{@{}p{0.08\textwidth}p{0.33\textwidth}p{0.09\textwidth}p{0.40\textwidth}@{}}
\toprule
ID & Formula or family & Status & Constructive proof location or qualification\\
\midrule\endhead
W1 & Integral definition on $\Rea a,\Rea b>0$ & \statusA & \Cref{sec:definition}; explicit endpoint and mesh modulus.\\
W2 & Symmetry $B(a,b)=B(b,a)$ & \statusA & Mesh reversal, \cref{thm:basic-laws}.\\
W3 & Beta--Gamma quotient & \statusB & One-dimensional effective limit, \cref{thm:beta-gamma}.\\
W4 & Product-of-integrals/Fubini proof & \statusC & Avoided; finite two-dimensional version is possible but unnecessary.\\
W5 & Convolution proof of beta--gamma & \statusB & One-dimensional convolution \eqref{eq:convolution} plus Gamma integration.\\
W6 & Positive integer factorial value & \statusA & \Cref{thm:integer-values}.\\
W7 & Binomial form of integer value & \statusA & Exact factorial algebra.\\
W8 & First derivative via digamma & \statusB & \Cref{thm:digamma-beta}.\\
W9 & Multivariate derivative & \statusB & \eqref{eq:multibeta-derivative}.\\
W10 & Two-parameter Stirling asymptotic & \statusB & Explicit relative error, \eqref{eq:beta-stirling-relative}.\\
W11 & Fixed-$y$ asymptotic & \statusB & Natural-grid modulus, \cref{thm:fixed-y-asymp}.\\
W12 & Trigonometric integral & \statusA & \eqref{eq:trig-beta}.\\
W13 & Infinite-interval integral & \statusA & \eqref{eq:infinity-beta}.\\
W14 & Power-substitution integral & \statusA & \eqref{eq:power-beta}.\\
W15 & Fractional-linear integral & \statusA & \eqref{eq:mobius-beta}; parameter must satisfy $c<1$.\\
W16 & Central radical integral & \statusA & \eqref{eq:central-radical}; positive real $z$, valid also at $z=1$.\\
W17 & General infinite series & \statusA & \Cref{thm:beta-series}.\\
W18 & Displayed symmetric series & \statusD & False as rendered; counterexample \cref{warn:false-series}.\\
W19 & Infinite product & \statusB & \Cref{thm:beta-product}; effective $O(1/N)$ tail.\\
W20 & Pascal identity & \statusA & \Cref{thm:pascal}.\\
W21 & Two coordinate recurrences & \statusA & \Cref{cor:coordinate-recurrence}.\\
W22 & $h(a,b)$ derivative values & \statusA & \Cref{thm:h-derivatives}.\\
W23 & PDE $h=h_a+h_b$ & \statusA & Same theorem, pointwise $t+(1-t)=1$.\\
W24 & Truncated-power convolution & \statusA & \eqref{eq:convolution}.\\
W25 & $B(1,x)=1/x$ & \statusA & Base case of \cref{thm:integer-values}.\\
W26 & Reflection $B(x,1-x)$ & \statusB & \Cref{cor:beta-reflection}; strip or meromorphic margin.\\
W27 & Product-reflection identity & \statusB & \Cref{cor:beta-reflection}.\\
W28 & Duplication identity & \statusB & \Cref{cor:beta-duplication}.\\
W29 & Pochhammer contour & \statusC & Polygon/branch route, \cref{sec:complex}.\\
W30 & Beta definition of $\binom nk$ & \statusA & \Cref{cor:binomial-beta}.\\
W31 & Continuous binomial interpolation & \statusB & Cleared reflection identity; raw quotient excludes integers.\\
W32 & Reciprocal-beta trig integrals & \statusA/\statusC & Integer frequencies finite; general frequency needs hypergeometric summation.\\
W33 & Incomplete beta definition & \statusA & \Cref{sec:incomplete}.\\
W34 & Incomplete trigonometric form & \statusA & \eqref{eq:inc-trig}.\\
W35 & Incomplete infinite-tail form & \statusA & \eqref{eq:inc-infinity}.\\
W36 & Hypergeometric representation & \statusA & \eqref{eq:inc-series}--\eqref{eq:inc-hyper}.\\
W37 & Regularized definition & \statusA & Guarded division using positivity.\\
W38 & Binomial CDF relation & \statusA & \Cref{thm:binomial-cdf}.\\
W39 & $I_0,I_1,I_x(a,1),I_x(1,b)$ & \statusA & \Cref{thm:incomplete-basic}.\\
W40 & Complement $I_x=1-I_{1-x}$ & \statusA & Split and reverse the finite integral.\\
W41 & Regularized recurrences & \statusA & \eqref{eq:inc-left-rec}--\eqref{eq:inc-right-rec}.\\
W42 & Endpoint derivative of $I_x$ & \statusA & Finite fundamental theorem.\\
W43 & Antiderivatives of $I_x$ and $B_x$ & \statusA & Product-rule telescope.\\
W44 & Branch transformation with $(-1)^a$ & \statusD/\statusC & Integer $a$ finite; complex version needs a branch.\\
W45 & $B_x(a,-a)$ & \statusA & Differentiate explicit power, \eqref{eq:inc-special}.\\
W46 & Continued fraction coefficients & \statusA & Primitive-recursive rational sequence.\\
W47 & Continuant determinant & \statusA & Second-order induction.\\
W48 & Alternating CF bounds & \statusC & Positive-remainder theorem remains.\\
W49 & Multivariate Gamma quotient & \statusB & Finite Gamma product.\\
W50 & Multivariate Pascal identity & \statusA/\statusB & \Cref{thm:multi-recurrence}.\\
M1 & Gamma/factorial definition & \statusA/\statusB & Integer case \statusA; general quotient \statusB.\\
{\scriptsize M2--M8} & Polar-coordinate factorial derivation & \statusC & Result proved one-dimensionally; finite polar-grid implementation omitted.\\
{\scriptsize M9--M10} & Odd-power trig integral and factorial value & \statusA & \Cref{thm:transformed-integrals,thm:integer-values}.\\
{\scriptsize M13--M19} & Symmetry and ordinary beta integral & \statusA & \Cref{thm:basic-laws,sec:definition}.\\
{\scriptsize M20--M21} & $u=x^2$ form & \statusA & Power substitution with $q=2$.\\
M22 & Infinite-interval polynomial form & \statusA & \eqref{eq:infinity-beta}.\\
{\scriptsize M23--M24} & $\int_0^\infty(1+x^p)^{-1}$ & \statusB & \eqref{eq:p-integral}.\\
{\scriptsize M25--M28} & First and second derivatives & \statusB & \Cref{thm:digamma-beta}.\\
{\scriptsize M29--M30} & General multiplication identity & \statusB & \Cref{thm:beta-multiplication}.\\
{\scriptsize M31--M35} & Shift, Pascal, recurrence & \statusA & \Cref{thm:pascal,cor:coordinate-recurrence}.\\
M36 & $B(p,n+1)$ finite product & \statusA & \Cref{thm:integer-values}.\\
M37 & Product of two central Betas & \statusB & \Cref{cor:beta-duplication}.\\
M38 & Associativity identity & \statusB & \Cref{cor:beta-assoc}.\\
M39 & Euler infinite product & \statusB & \Cref{thm:beta-product}.\\
M40 & Gosper odd product & \statusB & \Cref{thm:gosper-odd}.\\
M41 & Gosper even product & \statusB & \Cref{thm:gosper-even}.\\
M42 & Three-factor special case & \statusB & $n=1$ in odd Gosper formula.\\
M43 & Four-factor special case & \statusB & $n=2$ in even Gosper formula.\\
\bottomrule
\end{longtable}

The ledger deliberately distinguishes ``identity proved from an effective Gamma contract'' from ``all Gamma obligations already internalized.''  This is the main proof-status boundary in the packet.

\section{Further research questions}
\label{sec:further}

\begin{enumerate}
\item \textbf{A direct beta library independent of Gamma.}
Can reflection, duplication, and the full multiplication theory be derived from the beta integral and finite trigonometric factorization without passing through the effective Gamma packet?  This may shorten the object-theory dependency graph.

\item \textbf{A sharp primitive-recursive beta evaluator.}
The evaluator in \cref{sec:definition} is intentionally wasteful.  Develop interval arithmetic using tanh--sinh or Gauss--Jacobi quadrature with a proof whose checker is still small and whose moduli are elementary enough for iPRA.

\item \textbf{Correct the symmetric-series source formula.}
Determine the intended identity behind the false Wikipedia display, locate its primary source, and produce a corrected Abel-summation theorem with explicit convergence hypotheses.

\item \textbf{Complete the incomplete-beta continued fraction.}
Construct the positive integral remainder for each continuant and prove the $4m$/$4m+1$ lower and $4m+2$/$4m+3$ upper pattern in \code{Prf PRAAx}.  This would provide an efficient certified evaluator for beta-distribution tails.

\item \textbf{Formalize Gauss hypergeometric summation.}
A single effective ${}_2F_1(1)$ theorem would complete the general reciprocal-beta trigonometric integrals, simplify the continued-fraction proof, and provide a reusable bridge to many other special functions.

\item \textbf{Branch-certified contour calculus.}
Build a generic polygonal-contour library where a branch of a logarithm is positive data.  Use it to internalize the Pochhammer contour and compare the contour evaluator with the Gamma quotient on overlapping domains.

\item \textbf{Variable-dimensional simplex integration.}
The fixed-dimensional multivariate beta proof is straightforward.  Can a single coded theorem quantify over the dimension while retaining only quantifier-free induction invariants and avoiding a universal evaluator?

\item \textbf{Probability without measure theory.}
Develop the beta and Dirichlet distributions as finite weighted-sum names with explicit normalization and moment certificates.  Prove concentration and posterior-update formulas directly in iPRA.

\item \textbf{Optimal proof-checker size.}
Among the integral, Euler-product, and Gamma-quotient definitions of Beta, which yields the smallest complete certificate checker after all necessary arithmetic and elementary-function code is included?

\item \textbf{Automated source auditing.}
Extend the exact-rational checker so that a web formula is translated to a symbolic rising/falling-factorial convention and tested automatically at small rational parameters before it enters a proof packet.
\end{enumerate}

The most immediate completion target is the chain
\[
 \text{finite integral}
 \longrightarrow
 \text{recurrences}
 \longrightarrow
 \text{effective beta--gamma limit}
 \longrightarrow
 \text{object-theory derivations}.
\]
It would certify the majority of the source-page formulas without first building the contour or continued-fraction layers.


\begin{thebibliography}{99}
\bibitem{wiki} Wikipedia contributors, ``Beta function,'' \url{https://en.wikipedia.org/wiki/Beta_function}, accessed 16 November 2026.
\bibitem{mathworld} E. W. Weisstein, ``Beta Function,'' MathWorld--A Wolfram Web Resource, \url{https://mathworld.wolfram.com/BetaFunction.html}, accessed 16 November 2026.
\bibitem{dlmf} NIST Digital Library of Mathematical Functions, Chapters 5 and 8, \url{https://dlmf.nist.gov/}, accessed 16 November 2026.
\bibitem{proveit} V. Reshetnikov, \emph{ProveIt}, \url{https://github.com/VladimirReshetnikov/ProveIt}, inspected on the main branch 16 November 2026.
\bibitem{gamma} ProveIt, \code{Logic/PrimitiveRecursiveArithmetic/docs/proofs/analysis/gamma_ipra}, effective Gamma research packet, inspected 16 November 2026.
\bibitem{concise} ProveIt, \code{Logic/PrimitiveRecursiveArithmetic/ConciseIPRA}, Lean formalization of theories bi-interpretable with iPRA.
\bibitem{derivations} ProveIt, \code{Logic/PrimitiveRecursiveArithmetic/Derivations}, kernel-checked embedded iPRA arithmetic library.
\end{thebibliography}

\end{document}
